\documentclass[11pt,letterpaper]{article}
\usepackage[T1]{fontenc}
\usepackage[utf8]{inputenc}
\usepackage{amsmath,amsthm}
\usepackage{newtxtext,newtxmath}
\usepackage{microtype,mathtools}
\usepackage[margin=1in,footskip=28pt]{geometry}
\usepackage{booktabs,array,tabularx,enumitem,graphicx}
\usepackage{needspace,float}
\usepackage[round,authoryear]{natbib}
\usepackage{xcolor}
\usepackage{hyperref}
\definecolor{linkblue}{RGB}{28,53,103}
\hypersetup{colorlinks=true,linkcolor=linkblue,citecolor=linkblue,urlcolor=linkblue,
 pdftitle={Other Findings},
 pdfsubject={Representation, minimax mistakes, metric generation, contrastive learning, and randomization},
 bookmarksnumbered=true}
\setlength{\parindent}{1em}
\setlength{\parskip}{2pt plus .5pt minus .5pt}
\setlength{\emergencystretch}{2em}
\setlist{itemsep=2pt,topsep=4pt,leftmargin=1.7em}
\renewcommand{\arraystretch}{1.17}
\theoremstyle{plain}
\newtheorem{theorem}{Theorem}[section]
\newtheorem{lemma}[theorem]{Lemma}
\newtheorem{proposition}[theorem]{Proposition}
\newtheorem{corollary}[theorem]{Corollary}
\theoremstyle{definition}
\newtheorem{definition}[theorem]{Definition}
\newtheorem{example}[theorem]{Example}
\newtheorem{remark}[theorem]{Remark}
\numberwithin{equation}{section}
\newcommand{\N}{\mathbb N}
\newcommand{\R}{\mathbb R}
\newcommand{\eps}{\varepsilon}
\DeclareMathOperator{\dist}{dist}
\DeclareMathOperator{\supp}{supp}
\DeclareMathOperator{\diam}{diam}
\newcolumntype{Y}{>{\raggedright\arraybackslash}X}
\pagestyle{plain}
\setcounter{tocdepth}{1}
\begin{document}
\section*{Other Findings}
\tableofcontents
\vspace{12pt}
\noindent These notes record five additional theoretical findings developed from the language-generation library. Each finding gives the motivation, source papers, result, and proof. The comparison tables collect their relationships and assumptions; the mathematical audit corrections are incorporated throughout.

\smallskip
\noindent\textit{Notation.} We use $\N=\{0,1,2,\ldots\}$, $[n]=\{1,\ldots,n\}$, and $x_+=\max\{x,0\}$. In the representative-generation sections, $\Delta(A)$ denotes countably additive probability laws on $A$. Each finding specifies its own observation and success conventions.
\clearpage
% Sources: LGITL_next_research_2026_09_19/source_context/overlap_characterization.tex;
% research_2026_09_19/proofs/two_block_bottleneck.tex and its independent audit.
% Written-proof chapter; no claim of a complete Lean certificate.
\section{Fresh representation: the exact obstruction and a sharp two-block bound}
\label{ov:chapter}

\subsection{Motivation and source problem}
Representative language generation asks for new valid examples whose group proportions resemble those of the observed examples \citep{P09}. Overlapping groups create a difficulty that disjoint groups conceal: preserving each marginal separately may require placing mass in an intersection that the observations have already exhausted. The question here is when this is the only kind of obstruction, and how quickly its effect vanishes with sample size.

This finding has two parts. First, for a test family of finite dual VC dimension, uniformly accurate fresh approximation is possible exactly when every fixed number of completed tests has uniformly bounded finite cells. Second, when the tests split into two families of dual VC dimension at most one, an exact Boolean bottleneck formula gives the optimal coefficient $1/\sqrt2$ in the square-root rate under linear cell growth. The second hypothesis is stronger than merely having dual VC dimension two.

The source problem and group-matching objective come from \citet{P09}. The results below concern the existence of fresh probability laws on a known countably infinite target after arbitrary finite samples. They do not assume independent observations or claim an efficient algorithm for learning an unknown target. The findings are the particular characterizations and sharp bounds; their proofs use classical separation, VC approximation, and finite-tree rounding tools.

\subsection{Model and statements}
Let $L$ be countably infinite and let $\mathcal A$ be a family of subsets of $L$, called tests. We may adjoin the empty test without changing any quantity below. Its \emph{pointwise completion} is
\[
 \mathcal K=\overline{\mathcal A}\subseteq\{0,1\}^{L}.
\]
This is closure of individual tests, not Boolean closure. The space $\mathcal K$ is compact and metrizable. The dual VC dimension is the VC dimension of the point profiles: a finite collection of tests is shattered if every binary membership pattern on it is realized by a point of $L$.

A completed cell of arity at most $r$ is a set
\[
 C(J,\sigma)=\{x\in L:\mathbf1_A(x)=\sigma_A\text{ for every }A\in J\},
 \qquad J\subseteq\mathcal K,\quad |J|\le r.
\]
Define its finite-cell capacity by
\[
 \overline B_r=\sup\{|C(J,\sigma)|:0<|C(J,\sigma)|<\infty,\ |J|\le r\},
 \qquad r\ge1,
\]
with empty supremum zero. The value may be infinite.

For a finite forbidden set $F$ and a probability law $\mu$ supported on $F$, put
\[
 e_F(\mu)=\inf_{Q\in\Delta(L\setminus F)}
              \sup_{A\in\mathcal A}|\mu(A)-Q(A)|.
\]
All probabilities are countably additive. For a finite nonempty sample $S$, let $P_S$ be its uniform law and define
\[
 e_L(S)=e_S(P_S),\qquad
 E_L(n)=\sup_{\substack{S\subset L\text{ finite}\\|S|\ge n}}e_L(S).
\]
Freshness excludes all of $F$, including any additional forbidden points of zero $\mu$-mass. The definition uses an infimum and does not presume an optimizer exists.

\begin{theorem}[Completed-cell characterization]\label{ov:criterion}
If the dual VC dimension of $\mathcal A$ is finite, then
\[
 E_L(n)\longrightarrow0
 \quad\Longleftrightarrow\quad
 \overline B_r<\infty\text{ for every fixed integer }r\ge1.
\]
More quantitatively, write $v$ for a finite upper bound on the dual VC dimension and $\Pi_v(r)=\sum_{j=0}^{\min(v,r)}\binom rj$. For every $0<\eta<1$ there is an integer $m=m(v,\eta)$, independent of the target and its probability laws, such that, with $r=2m$,
\begin{equation}\label{ov:quantitative}
 E_L(n)\le 2\eta+\frac{\Pi_v(r)\overline B_r}{n}.
\end{equation}
\end{theorem}

\begin{theorem}[Sharp square-root bound]\label{ov:sharp}
Suppose $\mathcal A=\mathcal A_1\cup\mathcal A_2$, each block has dual VC dimension at most one, and
\[
 \overline B_r\le C_0r\qquad(r\ge1),\qquad C_0>0.
\]
Then
\begin{equation}\label{ov:sqrt}
 E_L(n)\le\sqrt{\frac{C_0}{2n}}\qquad(n\ge2C_0).
\end{equation}
The coefficient $1/\sqrt2$ is optimal, even in one fixed countable target with $C_0=1$ and with the tests in each block pairwise disjoint.
\end{theorem}

The proofs follow in full. The quantitative bound in Theorem~\ref{ov:criterion} is a general-purpose bound, whereas Theorem~\ref{ov:sharp} uses the special two-block structure to obtain the sharp rate and constant.

\subsection{Completion and finite signed witnesses}
\begin{lemma}[Completion preserves the problem]\label{ov:completion}
Completion preserves dual VC dimension and, for any two countably additive probabilities $P,Q$ on $L$,
\[
 \sup_{A\in\mathcal A}|P(A)-Q(A)|
 =\sup_{A\in\mathcal K}|P(A)-Q(A)|.
\]
Restricting the infimum defining $e_F(\mu)$ to finite-support fresh laws does not change its value.
\end{lemma}
\begin{proof}
Enumerate $L$. For $A\in\mathcal K$, choose original tests agreeing with $A$ on the first $j$ points. Their indicators converge pointwise to $\mathbf1_A$. Dominated convergence for $P$ and $Q$ proves the equality of suprema. If completed tests shatter a finite set of test coordinates, choose one point for each pattern and approximate each test on that finite witness set. The original tests then realize the same patterns, proving preservation of dual dimension. Finally, truncate a countable fresh law and move its remaining mass to a fixed fresh point. Every test probability changes by at most that remaining mass, which tends to zero.
\end{proof}

For $x\in L$, define the continuous evaluation $f_x(A)=\mathbf1_A(x)$ on $\mathcal K$. The probability profile $f_\mu(A)=\mu(A)$ is continuous: it is the uniformly convergent sum $\sum_x\mu\{x\}f_x$. For a potential $h:L\to\mathbb R$, write
\[
 G_{F,\mu}(h)=\mathbb E_\mu h-\sup_{y\notin F}h(y).
\]

\begin{lemma}[Separation]\label{ov:separation}
The discrepancy has the representation
\begin{equation}\label{ov:dual}
 e_F(\mu)=\sup_{\|\lambda\|_{\mathrm{TV}}\le1}
 G_{F,\mu}\left(x\mapsto\int_{\mathcal K}f_x\,d\lambda\right),
\end{equation}
where $\lambda$ ranges over finite signed regular Borel measures on $\mathcal K$.
\end{lemma}
\begin{proof}
In the Banach space $C(\mathcal K)$ with its supremum norm, let $H$ be the norm-closed convex hull of $\{f_y:y\notin F\}$. Finite convex combinations are exactly the profiles of finite-support fresh laws. Lemma~\ref{ov:completion} shows that their distance from $f_\mu$ is $e_F(\mu)$. The Hahn--Banach distance formula for a closed convex set expresses this distance as the supremum of $\ell(f_\mu)-\sup_{g\in H}\ell(g)$ over continuous linear functionals of norm at most one. The Riesz representation theorem identifies those functionals with the signed regular measures in the display, with functional norm equal to total variation. A linear functional has the same supremum on a set and its closed convex hull, so its supremum on $H$ is the fresh-point supremum in~\eqref{ov:dual}.
\end{proof}

We include the classical VC sampling argument~\citep{VC} needed to reduce these possibly non-atomic separators to finite tests. Signed-measure decomposition followed by dual-VC approximation is established machinery; Gurvits and Koiran use it for finite approximation of convex superpositions~\citep{GK}. No compactness of output laws is assumed.

\begin{lemma}[Uniform finite approximation]\label{ov:vc}
Let $\mathcal R$ be a countable measurable range family of VC dimension at most $v$ on a probability space with law $\pi$. For each $0<\eta<1$, an empirical law $\pi_m$ on $m=m(v,\eta)$ points satisfies
\[
 \sup_{R\in\mathcal R}|\pi_m(R)-\pi(R)|\le\eta.
\]
One may choose any sufficiently large integer $m$ satisfying
\[
 m\eta^2\ge2,\qquad 4\Pi_v(2m)e^{-m\eta^2/8}<1.
\]
\end{lemma}
\begin{proof}
The Sauer trace bound gives at most $\Pi_v(k)$ different traces on $k$ points. Its counting recurrence follows by separating traces according to the last coordinate: traces with both extensions have dimension at most $v-1$, while all prefixes have dimension at most $v$. Induction and the binomial recurrence give the bound, including the zero-coordinate and zero-dimension cases.

Draw $m$ independent points and an independent ghost sample of size $m$, with empirical laws $\pi_m,\pi'_m$. If the first sample has error greater than $\eta$ on some range, select the first such range in a fixed enumeration. Conditional Chebyshev, using variance at most $1/(4m)$, gives ghost error at most $\eta/2$ with probability at least $1/2$. Thus
\[
 \Pr\!\left(\sup_R|\pi_m(R)-\pi(R)|>\eta\right)
 \le2\Pr\!\left(\sup_R|\pi_m(R)-\pi'_m(R)|>\eta/2\right).
\]
Independently swap the two members of each pooled sample pair; this does not change their joint law. Conditional on the pooled points, a fixed trace gives an empirical difference $m^{-1}\sum_{j=1}^m\varepsilon_j a_j$, where the independent signs $\varepsilon_j$ are symmetric and $|a_j|\le1$. The exponential bound $\mathbb E e^{t\varepsilon_j a_j}\le e^{t^2a_j^2/2}$, followed by Markov's inequality with the optimal $t$, bounds its two-sided tail at $\eta/2$ by $2e^{-m\eta^2/8}$. There are at most $\Pi_v(2m)$ traces. A union bound therefore bounds the original bad-sample probability by $4\Pi_v(2m)e^{-m\eta^2/8}<1$. Some sample is good. The stated integers exist because the trace bound is polynomial in $m$ for fixed $v$ and the exponential factor decays.
\end{proof}

\begin{corollary}[Finite signed witnesses suffice]\label{ov:finite-witness}
If the dual VC dimension is at most $v<\infty$, every $\lambda$ in~\eqref{ov:dual} can, for each $0<\eta<1$, be replaced by a signed atomic measure supported on at most $2m(v,\eta)$ completed tests, with total variation at most one, such that its potential differs uniformly on $L$ by at most $\eta$. Consequently $e_F(\mu)$ is the supremum of $G_{F,\mu}(h)$ over finite potentials
\[
 h=\sum_{A\in J}c_A\mathbf1_A,\qquad
 J\subseteq\mathcal K\text{ finite},\quad\sum_{A\in J}|c_A|\le1.
\]
\end{corollary}
\begin{proof}
Normalize the nonzero positive and negative parts of the Jordan decomposition of $\lambda$. Apply Lemma~\ref{ov:vc} to each, with ranges $R_x=\{A\in\mathcal K:x\in A\}$. These ranges are clopen, indexed by countable $L$, and have VC dimension at most $v$. Multiply the two empirical measures by the original Jordan masses and subtract. Their total variation is at most the sum of those masses, at most one; their uniform potential error is at most $\eta$ times that sum. A uniform potential error $\eta$ changes both the $\mu$-average and the fresh supremum by at most $\eta$, so it changes $G_{F,\mu}$ by at most $2\eta$. Letting $\eta$ tend to zero in~\eqref{ov:dual} proves the assertion.
\end{proof}

\subsection{Proof of the completed-cell characterization}
\begin{proof}[Proof of Theorem~\ref{ov:criterion}]
Suppose first that $\overline B_r=\infty$ for some fixed $r\ge1$. There are arbitrarily large finite cells $C=C(J,\sigma)$ with $1\le |J|\le r$. For the sample $S=C$, every fresh point differs from $\sigma$ on at least one test. Hence every fresh law $Q$ satisfies
\[
 \sum_{A\in J}|Q(A)-P_C(A)|
 =\mathbb E_Q\sum_{A\in J}\mathbf1\{\mathbf1_A(x)\ne\sigma_A\}\ge1.
\]
Some completed test has discrepancy at least $1/r$; completion preserves discrepancy. Since such samples have arbitrarily large size, $E_L(n)\ge1/r$ for every $n$. Necessity does not require a VC hypothesis.

For sufficiency fix $\eta$, put $r=2m(v,\eta)$, and consider a sample $S$ of size $N$. By Corollary~\ref{ov:finite-witness}, any unit separator has gap at most $2\eta$ plus the gap of a finite potential $h$ using at most $r$ completed tests and coefficient norm at most one. Those tests partition $L$ into at most $\Pi_v(r)$ nonempty cells. Move the empirical mass in each infinite cell to a fresh point in that same cell. Move the empirical mass in all finite cells to any fresh point. This produces a finite-support fresh law $Q$. The total empirical mass moved between distinct cells is at most
\[
 \frac{\Pi_v(r)\overline B_r}{N}.
\]
Moving one unit of mass changes $h$ by at most $\sum_A|c_A|\le1$, while within-cell moves do not change it. Therefore
\[
 G_{S,P_S}(h)\le\mathbb E_{P_S}h-\mathbb E_Qh
 \le\frac{\Pi_v(r)\overline B_r}{N}.
\]
Take the supremum over separators and then over $N\ge n$ to obtain~\eqref{ov:quantitative}. If every $\overline B_r$ is finite, first choose $\eta$ small and then $n$ large. This proves convergence.
\end{proof}

\subsection{Boolean costs in a VC-one block}
For the remainder assume $\mathcal A=\mathcal A_1\cup\mathcal A_2$ with both blocks of dual VC dimension at most one. Put $\mathcal K_i=\overline{\mathcal A_i}$; then $\mathcal K=\mathcal K_1\cup\mathcal K_2$. The full family has dual VC dimension at most two, because a shattered collection cannot contain two tests assigned to the same block.

Let $\mathcal B_i$ be the algebra of finite Boolean combinations of tests in $\mathcal K_i$. For $U\in\mathcal B_i$, define
\begin{equation}\label{ov:cost}
 \gamma_i(U)=\inf\left\{\sum_{A\in J}|c_A|:
 \mathbf1_U=c_0+\sum_{A\in J}c_A\mathbf1_A,\quad
 J\subseteq\mathcal K_i\text{ finite}\right\}.
\end{equation}
The additive constant is free, since it cancels when comparing probabilities.

\begin{lemma}[Tree representation and integral cost]\label{ov:tree}
Every $U\in\mathcal B_i$ has an attained, nonnegative integer cost. It has cost zero exactly when $U=\varnothing$ or $U=L$. An attaining representation is given by a binary labeling of a finite inclusion tree, with cost equal to the number of edges across which the label changes.
\end{lemma}
\begin{proof}
Fix $z\in L$ and finitely many tests from one block. Complement each test containing $z$, discard empty sets, and identify duplicates. The resulting sets all exclude $z$ and are laminar: overlapping, nonnested sets would realize $11,10,01$, while $z$ realizes $00$, contradicting dual VC dimension one. Form their inclusion tree with root $L$, attaching each set to its smallest strict superset. The residual atom of a vertex is its set minus its immediate children. These atoms partition $L$; some are empty and remain as vertices.

For any node values $p_v$, the function taking value $p_v$ on the residual atom at $v$ is
\begin{equation}\label{ov:telescoping}
 p_{\mathrm{root}}+
 \sum_{v\ne\mathrm{root}}(p_v-p_{\mathrm{parent}(v)})\mathbf1_{B_v}.
\end{equation}
Indeed, at an actual point only the sets along its root-to-atom path contribute, and their differences telescope. Replacing oriented indicators by original tests only changes coefficient signs and the free constant. A Boolean combination is constant on each residual atom, so binary node values representing it exist. Its cost is the number of crossing edges.

Conversely, start with any real-coefficient representation of $\mathbf1_U$. Orientation and combining duplicates do not increase coefficient norm. Assign root-path sums to the nodes; their total edge variation is at most that norm. Clamp every node value to $[0,1]$. This preserves values on nonempty atoms and cannot increase variation. Thresholding these values at $t\in(0,1)$ gives a binary representation of the same set $U$. If $k(t)$ counts its crossing edges, then
\[
 \int_0^1 k(t)\,dt
 =\sum_{v\ne\mathrm{root}}|p_v-p_{\mathrm{parent}(v)}|,
\]
because an edge is cut for an interval of thresholds whose length is its endpoint difference. Some threshold has integer cost no larger than the original norm. The nonempty set of all obtainable integer costs has a minimum, and this minimum is no larger than every real representation cost. It is itself an allowable representation cost, proving equality with~\eqref{ov:cost} and attainment. A zero-edge binary labeling is constant on its tree. Conversely, the empty and full sets have free constant representations.
\end{proof}

\subsection{The exact exhausted-intersection bottleneck}
\begin{theorem}[Exact two-block discrepancy]\label{ov:bottleneck}
For every finite $F$ and probability $\mu$ supported on $F$,
\begin{equation}\label{ov:exact}
 e_F(\mu)=
 \sup_{\substack{U\in\mathcal B_1,\ V\in\mathcal B_2\\
 U\cap V\subseteq F,\ \gamma_1(U)+\gamma_2(V)>0}}
 \frac{(\mu(U)+\mu(V)-1)_+}{\gamma_1(U)+\gamma_2(V)}.
\end{equation}
The empty supremum is zero. If this supremum is positive, it is attained by a pair $U,V$.
\end{theorem}
The numerator is an unavoidable marginal deficit: fresh output cannot put total mass greater than one into two sets whose intersection has been exhausted. The denominator converts discrepancies of these Boolean sets into discrepancies of the original tests.

\begin{proof}
Denote the right side by $D$. A fresh law satisfies $Q(U)+Q(V)\le1$. If its test discrepancy is $\delta$, the attained representations in Lemma~\ref{ov:tree} give $|Q(U)-\mu(U)|\le\delta\gamma_1(U)$ and its analogue for $V$. Thus every displayed ratio is at most $\delta$. Taking infima gives $D\le e_F(\mu)$.

For the reverse inequality, use Corollary~\ref{ov:finite-witness}. Assign each test in a finite unit-norm potential to one block and write $h=f+g$. Construct the two inclusion trees as in Lemma~\ref{ov:tree}, with node potentials representing $f$ and $g$. Their combined edge variation is at most one. Set
\[
 u=\max_{x\in L}f(x),\quad v=\max_{x\in L}g(x),\quad
 M=\max_{y\notin F}h(y),\quad \Delta=u+v-M\ge0.
\]
These maxima exist because the potentials take finitely many values. If $\Delta=0$, the separator gap is nonpositive. Otherwise apply to every node, including those with empty atoms, the clamping maps
\[
 \alpha=\left[\frac{f-u+\Delta}{\Delta}\right]_0^1,
 \qquad
 \beta=\left[\frac{g-v+\Delta}{\Delta}\right]_0^1,
 \qquad [a]_0^1=\min\{1,\max\{0,a\}\}.
\]
The display specifies the functions on actual atoms as well. Their combined tree variation is at most $1/\Delta$. On actual points the unclamped expressions are at most one, so clamping only increases them. Consequently
\begin{equation}\label{ov:gap-clamp}
 G_{F,\mu}(h)\le\Delta(\mathbb E_\mu\alpha+\mathbb E_\mu\beta-1).
\end{equation}
For fresh $y$, one has $\alpha(y)+\beta(y)\le1$: if either unclamped value is nonpositive, its clamped value is zero and the other is at most one; if both are positive, their sum is at most one by $h(y)\le M$.

For $0<t<1$, let $U_t=\{\alpha>t\}$ and $V_t=\{\beta\ge1-t\}$. They lie in their respective Boolean algebras and $U_t\cap V_t\subseteq F$, since membership in both gives $\alpha+\beta>1$. Thus
\[
 \mu(U_t)+\mu(V_t)-1\le D\bigl(\gamma_1(U_t)+\gamma_2(V_t)\bigr).
\]
This also holds when both costs vanish: both sets are constant and cannot both be $L$, because $L$ is infinite and their intersection is contained in finite $F$. Thresholding the node potentials bounds each Boolean cost by its tree boundary count. The edge identity used in Lemma~\ref{ov:tree} therefore yields
\[
 \int_0^1\bigl(\gamma_1(U_t)+\gamma_2(V_t)\bigr)\,dt
 \le\frac1\Delta.
\]
All integrands take finitely many values; strict versus weak thresholds affect only finitely many endpoints. Integrating the preceding inequality and using
$\int_0^1\mathbf1_{\{\alpha>t\}}dt=\alpha$ and
$\int_0^1\mathbf1_{\{\beta\ge1-t\}}dt=\beta$
in~\eqref{ov:gap-clamp} gives $G_{F,\mu}(h)\le D$. Finite signed witnesses suffice, so $e_F(\mu)\le D$.

Finally, there are finitely many trace pairs $(U\cap F,V\cap F)$. Each fixes the numerator. For every attainable trace pair admitting a positive denominator, the possible denominators are positive integers and have an attained minimum. Taking the best ratio over these finitely many trace pairs proves positive bottleneck attainment. This is attainment of a dual witness; it is not by itself a proof that the primal fresh-law infimum is attained.
\end{proof}

\subsection{Boundary counting gives the sharp constant}
\begin{proof}[Proof of the upper bound in Theorem~\ref{ov:sharp}]
We first allow arbitrary $\mu$, set $\rho=\max_x\mu\{x\}$, and write $a=C_0\rho$. Take an admissible pair in~\eqref{ov:exact} with positive costs $p=\gamma_1(U)$ and $q=\gamma_2(V)$. Choose attaining binary tree representations. Remove crossing edges. Each component labeled one has at least one boundary edge, and their total boundary size is $p$; thus there are at most $p$ such components. The corresponding assertions hold with $q$ in the other tree.

The union of residual atoms in a connected component is a completed cell with at most its boundary size in literals. To see this, if its highest vertex is a nonroot set $A$, the cell is $A$ minus the sets at all downward exit edges. Its defining literals are the entering positive literal and the exit negative literals. If the component contains the root, omit the entering literal. Oriented sets are literals of original completed tests. This argument also handles empty atoms.

Intersect every selected component cell from the first tree with every selected component cell from the second. The nonempty intersections partition $U\cap V$ and lie in finite $F$. If their component boundary sizes are $b_j$ and $c_k$, their defining arities are at most $b_j+c_k$. Therefore
\[
 |U\cap V|\le C_0\sum_{j,k}(b_j+c_k)
 \le C_0(qp+pq)=2C_0pq.
\]
Empty intersections may be retained in this upper bound. The numerator in~\eqref{ov:exact} is at most both $1$ and $\mu(U\cap V)$. Since $pq\le(p+q)^2/4$, putting $s=p+q$ bounds the ratio by
\[
 \min\left\{\frac1s,\frac{2apq}{s}\right\}
 \le\min\left\{\frac1s,\frac{as}{2}\right\}
 \le\sqrt{a/2}.
\]
The last inequality follows by considering $s$ below or above $\sqrt{2/a}$.

If one cost is zero, only the case in which its set equals $L$ can contribute a positive numerator. The other set is then contained in $F$. Its selected component cells have total boundary equal to its positive cost, so their total size is at most $C_0$ times that cost. This branch's ratio is at most $a$. Consequently
\[
 e_F(\mu)\le\min\{1,\max\{a,\sqrt{a/2}\}\}.
\]
For the uniform law on a sample of size $N\ge n\ge2C_0$, $a=C_0/N\le1/2$, so $a\le\sqrt{a/2}$. Taking the supremum over those samples proves~\eqref{ov:sqrt}.
\end{proof}

\subsection{A single target witnessing optimality}
\begin{proof}[Sharpness in Theorem~\ref{ov:sharp}]
For every $k\ge1$, create $k$ infinite private row reservoirs $R_{k,i}$, $k$ infinite private column reservoirs $T_{k,j}$, and a two-point set $F_{k,ij}$ for each pair $1\le i,j\le k$. All these sets, including sets from different values of $k$, are disjoint. Add an infinite background. Their union is one countable target $L$. Define row and column tests
\[
 U_{k,i}=R_{k,i}\mathbin{\dot\cup}\bigcup_{j=1}^kF_{k,ij},
 \qquad
 V_{k,j}=T_{k,j}\mathbin{\dot\cup}\bigcup_{i=1}^kF_{k,ij}.
\]
Within each block the tests are pairwise disjoint, hence have dual VC dimension at most one. Each point belongs to at most two tests. Any sequence of distinct tests converges pointwise to the empty set, while a convergent sequence with a repeated test subsequence has that test as its limit. Thus completion adds only the empty test.

A nonempty cell with no positive nonconstant literal contains the infinite background. A cell with exactly one positive test retains that test's entire private reservoir, unless it is inconsistent. A cell with two compatible positive tests is one pair set $F_{k,ij}$; additional negative literals either preserve or empty it, because its two points have identical profiles. Two positives from the same block, or from incompatible gadgets, give the empty cell. It follows that every nonempty finite completed cell has size two and least arity two. Hence $\overline B_1=0$ and $\overline B_r=2$ for $r\ge2$, in particular $\overline B_r\le r$.

Take $S_k=\bigcup_{i,j}F_{k,ij}$, of size $2k^2$. Its total empirical mass summed over the $2k$ local tests is two. Every fresh point belongs to at most one of those tests, so every fresh law has total local test mass at most one. At least one test therefore has discrepancy at least $1/(2k)$. Put mass $1/(2k)$ on one point from each private row and column reservoir. Each local test then has mass $1/(2k)$, compared with empirical mass $1/k$; all other tests have mass zero under both laws. This fresh law attains error $1/(2k)$. Thus
\[
 e_L(S_k)=\frac1{2k}=\sqrt{\frac1{2|S_k|}}
\]
for every $k$ in the same target, proving sharpness. At $k=2$, eight observed points have row and column proportions $1/2$; every fresh distribution incurs error at least $1/4$, and a four-point reservoir law achieves it.
\end{proof}

\subsection{Interpretation and scope}
Theorem~\ref{ov:criterion} separates qualitative feasibility from quantitative rates. Finiteness of the capacity at each fixed arity is sufficient for convergence, whereas a prescribed growth bound yields an explicit rate. In the two-block setting, Theorem~\ref{ov:bottleneck} further reduces discrepancy to exhausted intersections of Boolean sets. Their boundary costs control the finite intersection and give the optimal square-root constant. The argument allows arbitrary correlations between the two blocks.

% Chapter 2. Adapted from the archived sharp_log_coefficient.tex (2026-09-18).
% Full proofs retained; concentration and unrelated extensions omitted.
\section{The exact logarithmic cost of representative generation}
\label{gm:chapter}
\newcommand{\gmip}[2]{\langle #1,#2\rangle}
\newcommand{\gmpos}[1]{\left(#1\right)_+}

\subsection{Motivation and source problems}
A representative generator should produce new valid points while preserving the proportions of groups in its observed data. These demands can conflict even when the target is known. Suppose a group contains exactly $B$ target points, all of which have already appeared. After $t\ge B$ observations its empirical share is $B/t$, but no fresh valid point can carry that share. Matching it within $\alpha$ forces at least $(B/t-\alpha)_+$ probability of an invalid or already observed output. Summing these obligations gives a logarithmic cost as the tolerance shrinks.

This chapter combines the representation objective of \emph{Representative Language Generation} by \citet{P09} with the cumulative-error question of \emph{Mistake-Bounded Language Generation} by \citet{P29}. The protocol change is explicit: the latter requires freshness at every round and counts invalid outputs; here representation is required at every round, and an output is charged if it is invalid \emph{or already observed}. Thus the temporary conflict is measured, and successful outputs still satisfy both validity and freshness. This is a combined model, not an unchanged application of either source characterization.

The result is an exact leading coefficient $\Gamma$ for the optimal total expected cost for each fixed finite target family and finite group family. It accounts jointly for exhausted group profiles and uncertainty about the true target. Two examples make the coefficient concrete: private target tails give $\Gamma=B(1-1/N)$, whereas parity profiles give $\Gamma=1$ despite exponentially many finite profile points. We also prove that deciding whether the coefficient is at least a given rational threshold is NP-complete for an explicit finite signature.

\subsection{The model and the main theorem}
Let $X$ be countably infinite, let $\mathcal H=(H_1,\ldots,H_N)$ be a finite nonempty list of infinite subsets of $X$, and let $\mathcal G=(G_1,\ldots,G_m)$ be a finite family of groups, with $m\ge1$. Points outside $\bigcup_iH_i$ are allowed in the universe. Write
\[
v(x)=\bigl(\mathbf 1\{x\in G_1\},\ldots,\mathbf 1\{x\in G_m\}\bigr)\in\{0,1\}^m.
\]
We allow the groups to leave points uncovered. If a covering family is required, as in \citet{P09}, append the universal group $X$; its constraint is automatic and its coordinate cancels from every profile displacement, leaving the minimax value and coefficient unchanged.

The target $H_z$ is fixed and unknown. At round $t\ge1$, the learner has seen distinct positive observations $x_1,\ldots,x_t\in H_z$ and chooses a probability law $Q_t$ on $X$. The next observation can be selected after the adversary sees the current draw. No correctness feedback is supplied. Put
\[
S_t=\{x_1,\ldots,x_t\},\qquad p_t=\frac1t\sum_{x\in S_t}v(x).
\]
For a fixed $0<\alpha<1$, representation requires
\begin{equation}\label{gm:eq:representation}
\left\|\mathbb E_{Q_t}v-p_t\right\|_\infty\le\alpha
\end{equation}
at every positive history. A draw $Y_t$ is a mistake when $Y_t\notin H_z\setminus S_t$. Thus previously generated values can be repeated if they have not been observed; freshness concerns the input sample.

Let $V(\alpha)$ be the infimum, over representative generators, of the supremum expected total number of mistakes over fixed targets and legal adaptive adversaries. The upper construction below is a deterministic function of positive history into finite-support laws. The lower bound also permits learners to retain previous random draws. We may require adversarial streams to be complete injective presentations, or allow arbitrary injective valid streams; all stated bounds hold under either convention.

For nonempty $F\subseteq[N]$, let $K_F=\bigcap_{i\in F}H_i$. Define
\[
B_{\mathrm{fin}}=\max\bigl(\{0\}\cup
\{|K_F|:F\ne\varnothing,\ K_F\text{ finite}\}\bigr).
\]
When $K_F$ is infinite, let $I_F$ be the profiles with infinite cells in $K_F$, and let $P_F$ be the profiles with nonempty finite cells. For $v\in P_F$ write $c_v=|\{x\in K_F:v(x)=v\}|$, with dependence on $F$ understood. Set
\[
B_{\mathrm{inf}}=\max_{F:K_F\text{ infinite}}\sum_{v\in P_F}c_v,
\qquad T_0=\max\{1,B_{\mathrm{fin}}+1,mB_{\mathrm{inf}}\}.
\]
All these numbers are finite. At least one infinite intersection exists because each individual target is infinite.

\begin{theorem}[Sharp logarithmic minimax law]\label{gm:thm:main}
There is an explicitly defined rational number $\Gamma=\Gamma(\mathcal H,\mathcal G)$, given in Definition~\ref{gm:def:gamma}, such that
\[
0\le\Gamma\le B_{\mathrm{inf}},\qquad
V(\alpha)=\Gamma\ln(1/\alpha)+O_{\mathcal H,\mathcal G}(1)
\quad\text{as }\alpha\downarrow0.
\]
For every $0<\alpha<1$, a representative generator satisfies
\begin{equation}\label{gm:eq:upper-main}
\sum_{t\ge1}\Pr\{Y_t\notin H_z\setminus S_t\mid\mathcal F_t\}
\le T_0-1+\Gamma\ln(1/\alpha)
\end{equation}
on every legal history, where $\mathcal F_t$ denotes information before the current draw. In particular the same expression bounds expected total mistakes.

If $\Gamma>0$, there are explicitly obtainable constants $A,W>0$ such that, whenever $0<\alpha<\Gamma/(AW)$,
\begin{equation}\label{gm:eq:lower-main}
V(\alpha)\ge\Gamma\ln(1/\alpha)
-\Gamma\left(1+\ln\frac{AW}{\Gamma}\right)+\alpha AW.
\end{equation}
Moreover $\Gamma=0$ exactly when every infinite target intersection has no nonempty finite profile cell.
\end{theorem}
The constant in the upper remainder is explicit. This matters because a general parametric linear-program argument would yield only an unspecified stabilization threshold. The binary profile structure provides the sharper bound through a simple absence-of-cancellation property.

\subsection{The polyhedral coefficient}\label{gm:sec:coefficient}
Fix $F\ne\varnothing$ with $K_F$ infinite. For a point $x\notin K_F$, let $J_F(x)=\{i\in F:x\in H_i\}$. Let $O_F$ be the finite set of realized pairs $(v(x),J_F(x))$ outside $K_F$. Every such $J_F(x)$ is a proper subset of $F$.

\begin{definition}[The coefficient]\label{gm:def:gamma}
For an anchor $u\in I_F$, a vector $a\in\R^m$ and a probability vector $\mu\in\Delta(F)$ are feasible when
\begin{align}
\gmip{a}{w-u}&\le0 &&(w\in I_F),\label{gm:eq:dual-infinite}\\
\gmip{a}{v-u}&\le1 &&(v\in P_F),\label{gm:eq:dual-finite}\\
\gmip{a}{w-u}&\le1-\sum_{i\in J}\mu_i &&((w,J)\in O_F).\label{gm:eq:dual-outside}
\end{align}
Set
\[
\Gamma_F=\max_{\substack{u\in I_F\\(a,\mu)\text{ feasible}}}
\sum_{v\in P_F}c_v\gmpos{\gmip{a}{v-u}},
\qquad
\Gamma=\max_{F:K_F\text{ infinite}}\Gamma_F.
\]
\end{definition}
The affine potential $\phi(x)=\gmip{a}{v(x)-u}$ is nonpositive on safe infinite cells. It is bounded by one on finite cells that might be exhausted, and by the $\mu$-average membership error outside the common core. Its positive empirical mass therefore certifies an unavoidable error cost.

\begin{lemma}[Attainment, computation, and positivity]\label{gm:lem:coefficient}
The maxima in Definition~\ref{gm:def:gamma} are attained, and $\Gamma$ is rational and computable from the realized joint group/target-membership profiles and their finite or infinite cardinalities. Also,
\[
0\le\Gamma_F\le\sum_{v\in P_F}c_v.
\]
If $v\in P_F$, then $\Gamma_F\ge c_v/|F|$. Consequently $\Gamma=0$ if and only if $B_{\mathrm{inf}}=0$.
\end{lemma}
\begin{proof}
Replace the positive-part sum by the finite maximum, over $E\subseteq P_F$, of
$\sum_{v\in E}c_v\gmip{a}{v-u}$.
For fixed $F,u,E$, maximizing this expression under \eqref{gm:eq:dual-infinite}--\eqref{gm:eq:dual-outside} and $\mu\in\Delta(F)$ is a rational linear program. It is feasible at $a=0$ and its objective is bounded above by $\sum_{v\in E}c_v$. Finite linear-program attainment and rationality give the claims after taking the finite maxima. This is a finite procedure given the exact signature; it is not a procedure for discovering that signature from arbitrary membership programs.

For positivity, fix a finite profile $v$ and let $b_j=2v_j-1$. This direction uniquely maximizes $\gmip b w$ at $w=v$ among all binary profiles, since
$\gmip b{v-w}$ equals their Hamming distance. Choose $u\in I_F$ maximizing $\gmip b u$ there and put $d=\gmip b{v-u}>0$. Use $a=b/(|F|d)$ and the uniform prior on $F$. Then $\phi(v)=1/|F|$, every profile has potential at most $1/|F|$, and every outside-core point has average error at least $1/|F|$. Infinite core profiles have nonpositive potential. The pair is feasible, contributing at least $c_v/|F|$.
\end{proof}

\subsection{A tangent program for one sample}\label{gm:sec:tangent}
Fix an infinite core $K_F$, a sample $S\subset K_F$, and an anchor $u\in I_F$. Let $s_v$ be the number of sampled points in finite core profile $v$, and put
\[
b=\sum_{v\in P_F}s_v,\qquad
h=\sum_{v\in P_F}s_v(v-u).
\]
We use finitely many output types. An infinite or unexhausted finite core profile has a fresh representative with error vector zero. An exhausted finite profile has a stale representative with error vector one for every target in $F$. Finally, each outside-core type $(w,J)$ has error vector $d_i=\mathbf 1\{i\notin J\}$. Such points have never been observed because $S\subset K_F$.

Write the feature and error vectors of a type as $w_j$ and $d_j\in\{0,1\}^F$. Delete types with $w_j=u$: they have zero displacement and nonnegative costs, while a safe fresh anchor is available separately. Consider
\begin{equation}\label{gm:eq:primal}
\begin{aligned}
\text{minimize } &D\\
\text{subject to }&q_j\ge0,\qquad
\sum_jq_j(w_j-u)=h,\\
&\sum_jq_jd_{j,i}\le D\qquad(i\in F),
\end{aligned}
\end{equation}
where $D\in\R$. The constraints themselves imply $D\ge0$.

\begin{lemma}[Tangent value and mass]\label{gm:lem:tangent}
Program \eqref{gm:eq:primal} has an optimum with value at most $\Gamma_F$. Every feasible vector $q$ satisfies
\begin{equation}\label{gm:eq:mass-bound}
\sum_jq_j\le\|h\|_1\le mb.
\end{equation}
Moreover $\Gamma_F$ equals the maximum of these tangent values over all integer count vectors $0\le s_v\le c_v$ and anchors $u\in I_F$.
\end{lemma}
\begin{proof}
Feasibility is immediate: for each finite profile, send $s_v$ units to its fresh type if the cell is unexhausted, and otherwise to its stale type. Because $u$ is binary, every vector $w_j-u$ has the same coordinatewise signs: nonnegative where $u$ is zero, and nonpositive where $u$ is one. Thus
\[
\sum_jq_j\|w_j-u\|_1
=\left\|\sum_jq_j(w_j-u)\right\|_1=\|h\|_1.
\]
Every retained type has $\|w_j-u\|_1\ge1$, proving the first inequality in \eqref{gm:eq:mass-bound}; the second follows from $\|v-u\|_1\le m$. In particular feasible $q$ form a compact set, and the continuous objective $\max_i\sum_jq_jd_{j,i}$ attains its minimum.

For clarity, the dual of \eqref{gm:eq:primal} is
\begin{equation}\label{gm:eq:tangent-dual}
\begin{aligned}
\text{maximize }&\gmip a h\\
\text{over }&\mu\in\Delta(F),\ a\in\R^m,\\
\text{subject to }&\gmip a{w_j-u}\le\sum_i\mu_id_{j,i}\quad\text{for every type }j.
\end{aligned}
\end{equation}
Indeed, assigning multipliers $\mu_i\ge0$ to the error constraints and $a$ to the displacement equality gives Lagrangian
\[
D\Bigl(1-\sum_i\mu_i\Bigr)+\gmip a h
+\sum_jq_j\Bigl(\sum_i\mu_id_{j,i}-\gmip a{w_j-u}\Bigr).
\]
Its infimum over free $D$ and nonnegative $q$ is finite precisely under the displayed dual constraints. Finite linear-program duality applies.

Every dual-feasible pair satisfies Definition~\ref{gm:def:gamma}. A positive potential on a finite profile is possible only when that cell is exhausted, in which case $s_v=c_v$. Therefore
\[
\gmip a h=\sum_vs_v\gmip a{v-u}
\le\sum_vc_v\gmpos{\gmip a{v-u}}\le\Gamma_F.
\]
This proves the value bound. Conversely choose a maximizing pair for $\Gamma_F$ and exhaust exactly its positive-potential finite profiles, leaving the other finite counts zero. That pair is feasible for the resulting tangent dual, and its objective is $\Gamma_F$. Hence equality is attained for some count vector.
\end{proof}
The mass bound is the essential geometric step. It turns an infinitesimal replacement into an actual probability law after at most $mb$ observations. No limiting approximation or unproved algorithmic certificate is involved.

\subsection{Proof of the upper bound}
Before time $T_0$, use the empirical law. It is exactly representative and costs at most one mistake per round. At $t\ge T_0$, let $F$ be the version space of targets containing $S_t$. It is nonempty and its core is infinite, since it contains $t>B_{\mathrm{fin}}$ distinct points. Its finite sample count is $b\le B_{\mathrm{inf}}$.

For each anchor $u\in I_F$, take an optimal tangent flow $q^u$. Since its total mass is at most $mb\le t$, assign mass $q_j^u/t$ to a realizing point of each type, and put the remaining mass on a fresh anchor point. Call this law $Q_t^u$. Its feature mean is
\begin{equation}\label{gm:eq:anchor-mean}
u+\frac ht=
\frac{\sum_vs_vv+(t-b)u}{t},
\end{equation}
and its error probability is at most $\Gamma_F/t$ for every target in $F$.

If $t>b$, mix these laws with weights $n_u/(t-b)$, where $n_u$ counts observations in infinite core profile $u$. The resulting law $R_t$ matches $p_t$ exactly. When $t=b$, expression \eqref{gm:eq:anchor-mean} is independent of the anchor, so choose any one. In either case every consistent target has error probability at most $\Gamma/t$.

Let $A_t$ be the corresponding mixture of fresh anchor points. It is valid and fresh for every consistent target. Its mean differs from $p_t$ by at most $b/t\le B_{\mathrm{inf}}/t$ in each coordinate: only the finite-profile portion of the sample has been replaced. If $B_{\mathrm{inf}}=0$, the construction is already exactly representative and error-free from $T_0$ onward. Otherwise set
\[
\lambda_t=\min\{1,\alpha t/B_{\mathrm{inf}}\},\qquad
Q_t=(1-\lambda_t)R_t+\lambda_tA_t.
\]
The feature change is at most $\alpha$, and the error probability for every consistent target is at most
\begin{equation}\label{gm:eq:pointwise-error}
\Gamma\gmpos{\frac1t-\frac\alpha{B_{\mathrm{inf}}}}.
\end{equation}
This is a pointwise guarantee at every history, so it survives adaptive choices of later observations.

Put $B=B_{\mathrm{inf}}>0$. The positive summand in \eqref{gm:eq:pointwise-error} decreases in $t$. Since $T_0\ge B$, its sum from $T_0$ is at most its sum from $B$, which is bounded by its first value plus the integral from $B$ to $B/\alpha$:
\begin{align*}
\sum_{t\ge T_0}\Gamma\gmpos{1/t-\alpha/B}
&\le\Gamma\left[\ln(1/\alpha)-(1-\alpha)(1-1/B)\right]\\
&\le\Gamma\ln(1/\alpha).
\end{align*}
Adding the first $T_0-1$ rounds proves \eqref{gm:eq:upper-main}. Empirical fallback extends the law to histories outside the legal protocol.

\subsection{Proof of the matching lower bound}
Assume $\Gamma>0$. Choose an attaining $F,u,a,\mu$, and let
\[
E=\{v\in P_F:\gmip a{v-u}>0\},\qquad
W=\sum_{v\in E}c_v,\qquad A=\|a\|_1.
\]
Both $W$ and $A$ are positive. First present all $W$ points in these cells, then distinct points from the infinite core cell $u$. This prefix is legal for every target in $F$, and at each $t\ge W$ its empirical potential $\phi(x)=\gmip a{v(x)-u}$ is $\Gamma/t$.

For any output point $y$, its potential is bounded above by its $\mu$-average mistake indicator. An observed point has mistake indicator one and potential at most one. A fresh point inside $K_F$ has nonpositive potential, since every positive-potential finite cell has been exhausted. A point outside $K_F$ was not observed, and \eqref{gm:eq:dual-outside} bounds its potential by its average target-membership error. Consequently every representative output law has average mistake probability at least
\[
\Gamma/t-\alpha A.
\]
Along the common prefix, the learner receives the same observations under all targets in $F$, and no correctness feedback. Thus its output-law distribution is the same under all of them, even if the learner retains previous randomness. Averaging over the fixed prior $\mu$ lower-bounds the worst target's expected total mistakes.

Use the common prefix through $\lfloor\Gamma/(\alpha A)\rfloor$. For complete presentations, extend this finite prefix separately to an enumeration of each target. Those later observations cannot remove the mistakes already forced. For $\alpha<\Gamma/(AW)$, decreasing-sum comparison gives
\begin{align*}
V(\alpha)&\ge\sum_{t=W}^{\infty}\gmpos{\Gamma/t-\alpha A}\\
&\ge\int_W^{\Gamma/(\alpha A)}(\Gamma/x-\alpha A)\,dx\\
&=\Gamma\ln(1/\alpha)-\Gamma\left(1+\ln\frac{AW}{\Gamma}\right)+\alpha AW.
\end{align*}
This proves \eqref{gm:eq:lower-main}. Since every binary profile difference has infinity norm at most one, $AW\ge\Gamma$. Lemma~\ref{gm:lem:coefficient} handles the zero case and completes Theorem~\ref{gm:thm:main}.

\subsection{Two exact minimax families}
\subsubsection{Private target tails: the price of unresolved ambiguity}
Fix $N\ge2$ and $B\ge1$. Take pairwise disjoint sets $G,R,P_1,\ldots,P_N$, with $|G|=B$ and the others infinite. Let
\[
X=G\cup R\cup\bigcup_{i=1}^NP_i,\qquad
H_i=G\cup R\cup P_i,
\qquad G_1=G\cup\bigcup_{i=1}^NP_i.
\]
Each individual target has two infinite group-profile cells. If it were known, it would permit exactly representative fresh generation at every round. The obstruction arises only while the target is ambiguous.

\begin{proposition}[Exact private-tail value]\label{gm:prop:private}
For this one-group family and every $0<\alpha<1$,
\[
V(\alpha)=\left(1-\frac1N\right)
\sum_{t=B}^{\infty}\gmpos{B/t-\alpha},
\qquad \Gamma=B\left(1-\frac1N\right).
\]
\end{proposition}
\begin{proof}
Until a private point appears, all targets remain possible. If $G$ is not exhausted, each observed profile has a fresh common-core representative, so the learner can be exactly representative without errors. Once $G$ is exhausted, place group mass $(B/t-\alpha)_+$ uniformly on one fresh point from each of the $N$ private tails, and use a fresh point of $R$ for the remaining mass. This is representative and has mistake probability $(1-1/N)(B/t-\alpha)_+$ for every target. If a private point is observed, it identifies the target; its two infinite profile cells then permit exact fresh representation forever. The exhaustion time of $G$ is at least $B$, proving the upper bound.

For the lower bound, present $G$ first, then shared $R$ points through the positive part of the displayed series. Under a uniform target, each fresh group point has average error $1-1/N$, and each stale group point has error one. Representation forces group mass at least $(B/t-\alpha)_+$. Summing proves the lower bound; the common prefix can subsequently be completed to each target.

For a subfamily $F$ of size at least two, the core is $G\cup R$. Its anchor is zero, and the constraints on the scalar $a$ are $a\le1-\mu_i$ for $i\in F$. The optimum prior is uniform, giving $\Gamma_F=B(1-1/|F|)$. Singleton subfamilies have no finite profile cells. Maximizing gives the stated $\Gamma$.
\end{proof}
Thus the cost of preserving representation is not merely the sum of a target-identification term and a known-target penalty. Every individual target here has zero known-target penalty, yet their ambiguity creates a logarithmic cost with a precise dependence on $N$.

\subsubsection{Parity profiles: many exceptions, one logarithmic obstruction}
Fix $m\ge1$, and take a singleton class whose target equals $X$. Give every even-parity profile in $\{0,1\}^m$ an infinite cell and every odd-parity profile exactly one point. The groups test the coordinates.

\begin{proposition}[Exact parity value]\label{gm:prop:parity}
This family has $B_{\mathrm{inf}}=2^{m-1}$ but
\[
V(\alpha)=\sum_{t=1}^{\infty}\gmpos{1/t-m\alpha},\qquad \Gamma=1.
\]
For comparison, let $D_n$ be the worst, over samples of size $n$, minimum group discrepancy of a fresh law. Then
\[
D_n=\frac1{mn}\qquad(n\ge1).
\]
Thus the exact fresh-discrepancy coefficient $1/m$ is different from the exact cumulative-mistake coefficient one.
\end{proposition}
\begin{proof}
Before any odd point is observed, all sample mass can be moved within its infinite even-profile cell. Suppose exactly one odd point $v$ has been observed. Its empirical mass is $1/t$. Transfer up to $m\alpha$ of this mass uniformly to the $m$ even neighbors of $v$, which all have fresh representatives. Each coordinate changes by at most $\alpha$, leaving stale mass $(1/t-m\alpha)_+$. All other sample mass moves within infinite cells.

If at least two distinct odd points have been observed, their empirical profile average belongs to the convex hull of even profiles. Indeed, any two distinct odd vertices have a midpoint equal to the midpoint of two even vertices: flip one differing coordinate in both. Averaging these midpoint representations over all pairs reproduces the average of all observed odd vertices. All their empirical mass can therefore be realized on fresh even points without error. This proves the upper bound for every observation order.

For the lower bound, first observe an odd point $v$, then distinct points from one even neighbor through time $\lfloor1/(m\alpha)\rfloor$ when that time is positive. The affine potential
\[
\phi(w)=1-\operatorname{Ham}(v,w)
\]
is one at $v$, zero at its even neighbors, and nonpositive on every other fresh profile. Its empirical average is $1/t$, and the linear part has $\ell_1$ norm $m$. Thus representation forces mistake probability at least $(1/t-m\alpha)_+$. Afterwards complete the target presentation. If $m\alpha\ge1$, the upper bound is already zero, and the same exact formula holds. The case $m=1$ needs no two-odd-point case, since there is only one odd point; both remaining arguments apply unchanged.

For $D_n$, transferring the entire mass of a sole observed odd point to its even neighbors changes each coordinate by $1/(mn)$, while samples with zero or at least two odd points allow exact fresh matching. The same potential, with the sample consisting of one odd point and $n-1$ neighbor-profile points, forces discrepancy at least $1/(mn)$. Finally the exact mistake series has leading term $\ln(1/\alpha)$ for fixed $m$, so Theorem~\ref{gm:thm:main} gives $\Gamma=1$.
\end{proof}
This is an exponential separation between finite-profile counting and the logarithmic mistake coefficient. It also rules out simply replacing the count by the earlier discrepancy coefficient: instantaneous approximation and cumulative mistakes require different normalizations.

\subsection{Computational complexity of the coefficient}\label{gm:sec:complexity}

\subsubsection{The coefficient as a finite input problem}
Consider one known countably infinite target $X$ and finitely many binary groups, represented by a map $v:X\to\{0,1\}^m$. Suppose the profile-zero cell is infinite and every other realized profile has a finite positive cardinality. Write $\mathcal P$ for those finite profiles and $c_v$ for their cardinalities. The sharp logarithmic coefficient in this special case is
\begin{equation}\label{gm:eq:coefficient}
\Gamma(\mathcal P,c)=
\max_{a\in\R^m:\ a\cdot v\le1\ \forall v\in\mathcal P}
\sum_{v\in\mathcal P}c_v\gmpos{a\cdot v}.
\end{equation}
There are no outside-target points, and the only infinite-cell anchor is zero. The vector $a$ is unrestricted in sign. Negative coordinates are needed in the reduction below.

The input lists the binary profiles and their positive integer cardinalities, together with a rational threshold $K$. Absent profiles are understood to have empty cells. Every such list has a countably infinite realization: create $c_v$ distinct points at each listed profile and countably infinitely many zero-profile points. Thus the reduction concerns realizable group systems, rather than arbitrary infeasible signatures.

Our formulation permits the groups to leave some points uncovered. To impose the convention of \citet{P09} that the groups cover the universe, append the universal group $X$. Its representation constraint is automatic, and its coordinate cancels from every profile displacement, so both the minimax problem and $\Gamma$ are unchanged. The infinite anchor then has one nonzero coordinate. The covered version of the hardness result has finite profile Hamming weight at most four, and Hamming distance at most three from that infinite anchor.

\begin{theorem}\label{gm:thm:complexity}
Deciding whether $\Gamma(\mathcal P,c)\ge K$ is NP-complete. NP-hardness holds even when every finite profile has at most three nonzero entries and all cell cardinalities are bounded by a polynomial in the number of listed profiles. Consequently exact evaluation, or additive approximation with error strictly less than $1/2$, is NP-hard on these instances.

The threshold problem for the general finite-class coefficient is also in NP when its input is an explicit joint group/target-membership signature with finite cardinalities and infinite-cell flags. Hence that general threshold problem is NP-complete as well.
\end{theorem}
The finite cardinalities in the construction can be expanded into individual points with only polynomial growth. The hardness therefore does not arise from encoding an exponentially large finite cell by a short integer.

\subsubsection{An exact maximum-cut reduction}
Let $G=(\{1,\ldots,n\},E)$ be a finite simple undirected graph, write $e=|E|$, and set
\[
M=2e+1.
\]
Use $2n+3$ profile coordinates, labelled
\[
x_1,\ldots,x_n,\quad y_1,\ldots,y_n,\quad d,f,c.
\]
The same labels will denote the corresponding coordinates of the potential vector $a$. A set of coordinate labels denotes its binary incidence profile. Include the following finite cells:

\begin{table}[tbp]
\centering
\caption{Finite profile cells in the maximum-cut reduction. The zero-profile cell is infinite.}
\label{gm:tab:hardness}
\begin{tabular}{@{}lll@{}}
\toprule
Profiles & Number & Cardinality of each cell\\
\midrule
$\{d\},\ \{f\},\ \{c,d,f\}$ & $3$ & $M$\\
$\{x_i,y_i\}$ for $i\in[n]$ & $n$ & $M$\\
$\{x_i\},\ \{y_i\}$ for $i\in[n]$ & $2n$ & $1$\\
$\{x_i,y_j,c\},\ \{x_j,y_i,c\}$ for $\{i,j\}\in E$ & $2e$ & $1$\\
\bottomrule
\end{tabular}
\end{table}
Also include the infinite zero-profile cell. All listed finite profiles are distinct and have Hamming weight at most three. Call the first two rows the heavy profiles; there are $n+3$ of them.

\begin{lemma}[Exact identity]\label{gm:lem:identity}
For this instance,
\begin{equation}\label{gm:eq:identity}
\Gamma=M(n+3)+n+\operatorname{MaxCut}(G).
\end{equation}
\end{lemma}

\begin{proof}[Lower bound]
Take any cut, encoded by $x_i\in\{0,1\}$. Set
\[
y_i=1-x_i,\qquad d=f=1,\qquad c=-1.
\]
Each heavy profile has potential exactly one. Each singleton $x_i$ or $y_i$ has potential at most one. The two profiles associated with an edge have potentials $x_i-x_j$ and $x_j-x_i$, both at most one. Thus the potential is feasible in \eqref{gm:eq:coefficient}. Its heavy contribution is $M(n+3)$, its singleton contribution is $n$, and its two edge positive parts sum to $|x_i-x_j|$. A maximum cut proves the lower bound in \eqref{gm:eq:identity}.
\end{proof}

\begin{proof}[Upper bound]
Let $a$ be any feasible potential. Write $z_h$ for its value on heavy profile $h$, and define the total heavy deficit
\[
\sigma=\sum_{h\ \mathrm{heavy}}\bigl(1-(z_h)_+\bigr).
\]
Every $z_h\le1$, so each deficit lies in $[0,1]$. The heavy part of the objective is exactly $M(n+3)-M\sigma$.

The singleton constraints imply $x_i,y_i\le1$, and the heavy pair constraint implies $x_i+y_i\le1$. Hence
\begin{equation}\label{gm:eq:singleton-bound}
(x_i)_++(y_i)_+\le1.
\end{equation}
If both are nonnegative this follows from their sum; otherwise the positive one is at most one. Each oriented edge profile contributes at most one by feasibility. If $\sigma\ge1$, the full objective is therefore at most
\[
M(n+3)-M+n+2e
<M(n+3)+n\le M(n+3)+n+\operatorname{MaxCut}(G).
\]

It remains to consider $\sigma<1$. Every heavy profile then has strictly positive potential. We can write its deficits as
\begin{align*}
d&=1-\delta_d,& f&=1-\delta_f,& c+d+f&=1-\delta_0,\\
x_i+y_i&=1-\delta_i&&&&(i\in[n]),
\end{align*}
where all $\delta$'s are nonnegative and sum to $\sigma$. In particular
\[
c=-1+\delta_d+\delta_f-\delta_0.
\]
Set $\bar x_i=\max\{x_i,0\}$. Then $0\le\bar x_i\le1$, $x_i\le\bar x_i$, and
\begin{equation}\label{gm:eq:complement-bound}
y_i\le1-\bar x_i.
\end{equation}
Indeed, if $x_i\ge0$, then $y_i=1-\delta_i-x_i\le1-x_i$; if $x_i<0$, the singleton constraint $y_i\le1$ gives the claim. For every oriented edge,
\begin{align*}
x_i+y_j+c
&\le\bar x_i+(1-\bar x_j)-1+\delta_d+\delta_f\\
&=\bar x_i-\bar x_j+\delta_d+\delta_f.
\end{align*}
Since $(u+v)_+\le(u)_++v$ for $v\ge0$, the two orientations of an edge have total positive part at most
\[
|\bar x_i-\bar x_j|+2(\delta_d+\delta_f).
\]
For every vector $\bar x\in[0,1]^n$,
\begin{equation}\label{gm:eq:coarea}
\sum_{\{i,j\}\in E}|\bar x_i-\bar x_j|\le\operatorname{MaxCut}(G).
\end{equation}
One proof chooses a uniform threshold $U\in[0,1]$ and cuts at $\{i:\bar x_i\ge U\}$. The expected cut size is the left-hand side, and every cut has size at most $\operatorname{MaxCut}(G)$.

Combining \eqref{gm:eq:singleton-bound}--\eqref{gm:eq:coarea}, the objective is at most
\begin{align*}
M(n+3)-M\sigma+n+\operatorname{MaxCut}(G)+2e(\delta_d+\delta_f)
&\le M(n+3)+n+\operatorname{MaxCut}(G)+(2e-M)\sigma\\
&\le M(n+3)+n+\operatorname{MaxCut}(G).
\end{align*}
This proves the upper bound. It also shows that every maximizing potential has zero heavy deficit: the finite penalty forces all heavy profiles to have potential exactly one.
\end{proof}

\subsubsection{Completeness and the general signature}
Unweighted maximum-cut threshold decision is NP-complete~\citep{GJS1976}. Lemma~\ref{gm:lem:identity} gives the many-one reduction
\[
\operatorname{MaxCut}(G)\ge k
\quad\Longleftrightarrow\quad
\Gamma\ge M(n+3)+n+k.
\]
There are $3n+3+2e$ finite profiles, $2n+3$ coordinates, and each cell has cardinality at most $2e+1$. The total number of finite points is
\[
(2e+1)(n+3)+2n+2e,
\]
which is polynomial in the graph size. This proves the claimed restricted NP-hardness. The identity also shows that an additive approximation with error less than $1/2$ determines $\operatorname{MaxCut}(G)$ by rounding after subtraction of the known offset.

For membership in NP in the singleton case, guess a subset $\mathcal A\subseteq\mathcal P$ of profiles to contribute to the linear objective. The corresponding conditions are the rational linear inequalities
\[
a\cdot v\le1\quad(v\in\mathcal P),\qquad
\sum_{v\in\mathcal A}c_v(a\cdot v)\ge K.
\]
If these inequalities hold, the positive-part objective is at least the selected sum, so $\Gamma\ge K$. Conversely, an attaining potential with $\Gamma\ge K$ supplies $\mathcal A$ consisting of its positive-potential profiles. A feasible rational linear system has a rational feasible point of polynomial bit length, so the guessed subset and such a point form a polynomial certificate. Equivalently, an NP verifier can guess the subset and test this linear program in polynomial time.

For completeness, the same reasoning applies to the general coefficient. Its explicit input has one row for each realized pair of a group profile and a target-membership profile, marked by a finite integer cardinality or an infinite flag. A certificate first chooses a nonempty target subfamily $F$ and an infinite-cell anchor $u$. Scanning and grouping the input rows computes its infinite core profiles $I_F$, finite core profiles $P_F$ and their cardinalities $c_v$, and outside-core types $(w,J)$. It then guesses $\mathcal A\subseteq P_F$. The remaining feasibility problem is linear in a potential $a$ and a prior $\mu\in\Delta(F)$:
\begin{align*}
a\cdot(w-u)&\le0&&(w\in I_F),\\
a\cdot(v-u)&\le1&&(v\in P_F),\\
a\cdot(w-u)&\le1-\sum_{i\in J}\mu_i&&((w,J)\text{ an outside type}),\\
\sum_{v\in\mathcal A}c_v\,a\cdot(v-u)&\ge K.
\end{align*}
All dimensions, coefficient lengths, and grouping operations are polynomial in the explicit input size. Feasibility can again be certified or tested in polynomial time after the discrete guesses. The finite union of these feasible cases is exactly the event $\Gamma\ge K$. This proves membership in NP and completes Theorem~\ref{gm:thm:complexity}.

\subsubsection{Interpretation and limitations}
The logarithmic coefficient remains finitely determined and exactly characterized. Computational hardness enters when the number of groups grows: choosing which profile potentials should be positive already encodes a cut problem. The reduction already uses one known target, no outside-target outputs, and finitely many profiles. Even sparse profiles suffice.

The reduction rules out a general polynomial-time exact evaluator unless $\mathrm P=\mathrm{NP}$. It does not rule out efficient evaluation for fixed group dimension, special incidence structure, or useful approximation guarantees. Because the reduction has a large known additive offset, it does not by itself prove a constant-factor approximation barrier. Nor does coefficient hardness alone establish online running-time lower bounds for every generator that obtains a nonsharp mistake guarantee.



\subsection{Scope of the minimax characterization}
The asymptotic law fixes both the target family and the group family while $\alpha\downarrow0$; its additive constant need not be uniform in either family. The construction uses the exact joint profiles and their finite or infinite cardinalities. Consequently it is a semantic generator, rather than a polynomial membership-query algorithm. The complexity theorem concerns exact threshold evaluation of the sharp coefficient; it does not establish hardness of every useful generation guarantee or constant-factor inapproximability.

\section{Metric generation detects properness of the completion}
\label{mt:chapter}

\subsection{Motivation and source}
When examples live in a metric space, a new output can be different from every observation and still be arbitrarily close to something already seen. Metric language generation asks for a stronger form of novelty: the output must stay a prescribed positive distance from the whole observed sample. The amount of information in that sample is measured by a covering number. A thousand observations concentrated in one small ball may therefore count as less evidence than ten well-separated observations.

The covering-threshold framework comes from \emph{On Generation in Metric Spaces} by \citet{P21}. The finding developed here identifies a classical geometric property that the framework can detect. For a countable class of unbounded targets in a countable unbounded metric space, generation with a target-dependent threshold is possible at every positive pair of scales exactly when every nonempty finite intersection of targets is either unbounded or totally bounded. Consequently, the ability to generate for \emph{every} countable such class is equivalent to properness of the metric completion.

There is a concrete proof issue behind this connection. A ball cover may use centers outside the set it covers. Such centers cannot automatically be appended to a positive sample: to preserve every currently possible target, the new observations must belong to their common intersection. The library's diagnostic for this source paper isolates the distinction between ambient centers and internal centers. The proofs below work directly with legal positive samples and identify the obstruction that survives at all scales.

The all-scale qualification is essential. We allow the output radius to be larger than the input radius, and require success for \emph{every} positive pair. The result does not assert that each fixed-scale theorem in \citet{P21} has the same characterization. Unboundedness of every target also matters: it ensures that every fixed finite input threshold can actually be reached.

\subsection{The semantic model}

Fix a countably infinite set $X$ with an unbounded metric $m$.
A target is an unbounded subset of $X$, and a target class $\mathcal H$
is a collection of such subsets. A finite ordered history
$h=(x_1,\ldots,x_t)$ has observation set
$S(h)=\{x_1,\ldots,x_t\}$. Repetitions are permitted. It is positive for
$L$ when $S(h)\subseteq L$. A generator is a function from all finite
histories to $X$; no computability or oracle-efficiency assumption is made.

For $a>0$, write
\[
 B_m(x,a)=\{y\in X:m(x,y)\le a\},
 \qquad
 N_m(a;A)=\min\left\{|C|:
 \begin{array}{l}
 C\subseteq X\text{ finite},\\[-2pt]
 A\subseteq\displaystyle\bigcup_{c\in C}B_m(c,a)
 \end{array}\right\}.
\]
Set $N_m(a;A)=\infty$ when there is no finite cover, and
$N_m(a;\varnothing)=0$. Whenever finite covers exist, the minimum exists
because their cardinalities are natural numbers. Centers are ambient
points of $X$ and need not belong to $A$. A set is \emph{totally bounded}
if $N_m(a;A)<\infty$ for every $a>0$. It is \emph{bounded} if it is
contained in one ball of finite radius.

\begin{definition}[Uniform and nonuniform generation at all scales]
\label{mt:model}
For fixed input scale $a>0$ and output scale $b>0$, a generator $g$ has a
uniform threshold $D\ge1$ for $\mathcal H$ if, for every positive history,
\begin{equation}
\label{mt:guarantee}
 S(h)\subseteq L\in\mathcal H,\quad N_m(a;S(h))\ge D
 \quad\Longrightarrow\quad
 g(h)\in L\setminus\bigcup_{x\in S(h)}B_m(x,b).
\end{equation}
The class satisfies $\mathsf U_m$ if such a generator and threshold exist
for every $a,b>0$. It satisfies $\mathsf{NU}_m$ if, for every $a,b>0$,
there is one generator with a positive integer threshold $D_L$ for each
$L\in\mathcal H$, satisfying the same implication with $D_L$ in place
of $D$. Initially the generator may depend on $(a,b)$.
\end{definition}

Thus the nonuniform quantifiers are
\[
 \forall a,b>0\ \exists g\ \forall L\in\mathcal H\ \exists D_L\ge1
 \ \forall h:\quad \text{\eqref{mt:guarantee} holds with }D_L.
\]
The true target may influence the threshold, but it cannot select a different
generator. We will construct a single generator that in fact works for
all scale pairs, after changing only the thresholds. The empty class is
vacuous throughout; statements involving an enumeration will assume the
class is nonempty.

The model can equivalently be tested on finite prefixes of complete positive
presentations. Every positive finite history extends to such a presentation:
append an enumeration of its countable infinite target. Moreover, on any
complete presentation of an unbounded target, $N_m(a;S(h))$ eventually
crosses every finite threshold. To see this, select as many target points
as desired with pairwise distances greater than $2a$. At each selection
step, finitely many radius-$2a$ balls form a bounded set and cannot cover
the unbounded target. Once all selected points have appeared, every
radius-$a$ ball covers at most one of them.

\subsection{The geometric obstruction in a finite intersection}

\begin{lemma}[Packing and legal covering centers]
\label{mt:packing}
The following facts hold in any metric space.
\begin{enumerate}
\item If $N_m(a;A)=\infty$, then for every $q\ge1$ there are $q$ points
of $A$ whose pairwise distances are greater than $a$. A radius-$a/2$
cover of those points needs at least $q$ balls.
\item If $A$ is nonempty and totally bounded, then for every $b>0$ it
has a finite radius-$b$ cover whose centers belong to $A$.
\item If $A$ is unbounded and $S$ is finite, then for every $b>0$ some
point of $A$ has distance greater than $b$ from every point of $S$.
\end{enumerate}
\end{lemma}
\begin{proof}
For the first assertion, choose points greedily. If fewer than $q$ chosen
points already cover $A$ by their radius-$a$ balls, they give a finite
cover, contrary to the premise. Otherwise choose a point outside those
balls and continue. Two chosen points cannot belong to one radius-$a/2$
ball, since the triangle inequality would make their distance at most $a$.

For the second assertion, take a finite ambient radius-$b/2$ cover of $A$.
Discard every ball missing $A$, and choose a point of $A$ in each remaining
ball. If $y$ and the chosen point lie in the same original ball, their
distance is at most $b$. These chosen points are the required internal
centers. This argument changes the radius; it does not justify moving
ambient centers into $A$ at the original radius without a loss.

For the third assertion, a finite union of finite-radius balls is bounded.
For example, its distance from one fixed center is bounded by the largest
distance to the other finitely many centers plus $b$. It therefore cannot
contain $A$. A point outside that union has the required distances.
\end{proof}

For a nonempty finite family $\mathcal F\subseteq\mathcal H$, its common
intersection $A=\bigcap_{L\in\mathcal F}L$ is a region where positive
observations cannot distinguish the targets in $\mathcal F$.
If $A$ is bounded but not totally bounded, it can supply arbitrarily much
input-covering evidence while fitting inside a single sufficiently large
output-neighborhood. This is the obstruction.

\begin{lemma}[A bounded, non-totally-bounded intersection prevents generation]
\label{mt:obstruction}
If a nonempty finite subfamily $\mathcal F$ has a bounded intersection
$A$ that is not totally bounded, then $\mathcal H$ fails
$\mathsf{NU}_m$.
\end{lemma}
\begin{proof}
Choose $a>0$ with $N_m(a;A)=\infty$. The set $A$ is nonempty; fix
$c\in A$. Since $A$ is bounded, choose $b>0$ such that
$A\subseteq B_m(c,b)$.

Suppose a generator satisfies the nonuniform guarantee at input scale
$a/2$ and output scale $b$. Let $D$ be the maximum of its finitely many
thresholds $D_L$, $L\in\mathcal F$. By Lemma~\ref{mt:packing}, select
$D$ pairwise more-than-$a$ separated points of $A$. Form a history
containing these points and $c$. It is positive for every target in
$\mathcal F$, and its input cover number is at least $D$. Thus its one
output must belong to every target in $\mathcal F$, hence to $A$.
However, the observed point $c$ makes every point of $A$ non-novel at
radius $b$. This contradicts \eqref{mt:guarantee}.
The history extends to a complete presentation for each possible target,
so imposing that extension requirement does not remove the contradiction.
\end{proof}

\subsection{The exact finite-intersection criterion}

\begin{theorem}[Countable nonuniform generation]
\label{mt:intersection}
For a nonempty countable class $\mathcal H$ of unbounded targets,
the following are equivalent:
\begin{enumerate}
\item $\mathcal H$ satisfies $\mathsf{NU}_m$;
\item for every nonempty finite $\mathcal F\subseteq\mathcal H$, the
intersection $\bigcap_{L\in\mathcal F}L$ is either unbounded or totally
bounded.
\end{enumerate}
When these conditions hold, one generator serves every positive scale pair.
If $\mathcal H$ is finite, these conditions are also equivalent to
$\mathsf U_m$.
\end{theorem}
\begin{proof}
Necessity is Lemma~\ref{mt:obstruction}: an intersection that is neither
unbounded nor totally bounded is precisely a bounded, non-totally-bounded
set.

We first prove sufficiency for a finite class. For each realizable sample
$S$, define its version space and common core by
\[
 V(S)=\{L\in\mathcal H:S\subseteq L\},
 \qquad C(S)=\bigcap_{L\in V(S)}L.
\]
Here ``realizable'' means $V(S)\ne\varnothing$. Its core contains $S$.
There are only finitely many intersections of nonempty subfamilies of
$\mathcal H$. Under condition 2, every bounded one is totally bounded.
Fix an input radius $a>0$, and let $E_a$ be the maximum of their finite
covering numbers at radius $a$; set $E_a=0$ if there are no bounded
intersections.

If $N_m(a;S)>E_a$, the core $C(S)$ cannot be bounded, because a cover
of $C(S)$ also covers $S$. An unbounded core supplies a point outside
the radius-$b$ balls around any finite sample, by Lemma~\ref{mt:packing}.
Choosing such a point yields an output valid for every target consistent
with $S$. The threshold $E_a+1$ proves uniform generation at $(a,b)$.
For any history outside this guarantee, an arbitrary fixed output makes
the generator a total function. Uniform generation always implies
nonuniform generation by using the same threshold for every target.

For a countable class, enumerate its targets as $L_1,L_2,\ldots$,
repeating members if the class is finite. For a finite sample $S$, let
\[
 V_j(S)=\{L_i:1\le i\le j,\ S\subseteq L_i\},
 \qquad C_j(S)=\bigcap_{L\in V_j(S)}L
\]
whenever $V_j(S)$ is nonempty. At a history of length $t$, choose the
largest index $j\le t$ for which $V_j(S)$ is nonempty and $C_j(S)$ is
unbounded. If there is no such index, output a fixed fallback point.
Otherwise choose a point of $C_j(S)$ whose distance from each observation
is greater than $t$. Such a point exists by Lemma~\ref{mt:packing}.
This construction makes no reference to $a$ or $b$. Its choices may be
noncomputable, as allowed by the semantic model.

Fix a target whose first occurrence is $L_z$, and fix $a,b>0$.
Consider all intersections of nonempty subfamilies of
$\{L_1,\ldots,L_z\}$. Each bounded one is totally bounded, and there
are finitely many. Let $E_{z,a}$ be the maximum of their radius-$a$
covering numbers, with value zero when none are bounded. Choose a
positive integer
\[
 D_{L_z}>\max\{E_{z,a},z,b\}.
\]
Suppose a positive history for $L_z$ has $N_m(a;S)\ge D_{L_z}$.
Since a finite sample covers itself using its own points as centers,
\[
 N_m(a;S)\le |S|\le t.
\]
Thus $t>z$ and $t>b$. The version space $V_z(S)$ is nonempty because
it contains $L_z$. If its core were bounded, that core would have cover
number at most $E_{z,a}$, which would also bound the sample's cover
number. This contradicts the chosen threshold. Hence $z$ is eligible.

The selected index $j$ is consequently at least $z$. Since $L_z$
remains consistent and appears among the first $j$ targets,
$C_j(S)\subseteq L_z$. The output belongs to $L_z$ and has distance
greater than $t>b$ from every observation. This proves the required
target-dependent guarantee. The construction used one generator for all
targets and all scale pairs, completing the proof.
\end{proof}

\begin{remark}[What the scheduler needs]
The maximum in this construction is over the finite set
$\{1,\ldots,t\}$. The proof does not take a maximum of countably many
thresholds. For a fixed true target, it uses only finitely many earlier
targets and finitely many intersections. Countability of the class is the
resource that makes this finite-prefix argument available.
\end{remark}

\subsection{From a generation property to proper completion}

A metric space is \emph{proper} if every closed bounded subset is compact.
Let $\widehat X$ denote the metric completion of $X$, with $X$ identified
with its dense isometric copy. Completion is part of the conclusion:
the original space need not be complete or proper.

\begin{lemma}[Two targets realize every bounded obstruction]
\label{mt:tails}
If $X$ is unbounded and $A\subseteq X$ is bounded, there are disjoint
unbounded sets $T_0,T_1\subseteq X\setminus A$. Consequently, the
unbounded targets $L_0=A\cup T_0$ and $L_1=A\cup T_1$ have
intersection exactly $A$.
\end{lemma}
\begin{proof}
Fix $x_0\in X$. Recursively, for each $n\ge1$, choose two new points
$u_n,v_n$ outside $A$ and all points chosen earlier, with both distances
from $x_0$ greater than $n$. Make the two choices successively so that
they are distinct. Each choice excludes a bounded set: the union of $A$,
a finite set, and $B_m(x_0,n)$. Such a union cannot exhaust unbounded
$X$. Set $T_0=\{u_n:n\ge1\}$ and $T_1=\{v_n:n\ge1\}$.
They are disjoint, avoid $A$, and are unbounded by their distance
requirements. The asserted intersection follows immediately.
\end{proof}

\begin{theorem}[Generation characterizes proper completion]
\label{mt:proper}
For an unbounded countable metric space $(X,m)$, the following conditions
are equivalent:
\begin{enumerate}
\item Every countable class of unbounded targets satisfies
$\mathsf{NU}_m$.
\item Every two-target class of unbounded targets satisfies
$\mathsf U_m$.
\item Every bounded subset of $X$ is totally bounded.
\item The metric completion $\widehat X$ is proper.
\end{enumerate}
\end{theorem}
\begin{proof}
If condition 3 holds, every finite intersection of unbounded targets is
either unbounded or, by condition 3, totally bounded. Theorem~\ref{mt:intersection}
therefore gives conditions 1 and 2.

Conversely, if condition 3 fails, choose a bounded subset $A$ that is
not totally bounded. Lemma~\ref{mt:tails} supplies two unbounded targets
with intersection exactly $A$. Lemma~\ref{mt:obstruction} shows that
this two-target class fails even nonuniform generation at one positive
scale pair. It contradicts condition 1. Since uniform generation implies
nonuniform generation, it also contradicts condition 2. Thus either
condition 1 or condition 2 implies condition 3.

We next show that condition 4 implies condition 3. Let $A\subseteq X$
be bounded. Its closure in $\widehat X$ is closed and bounded, hence
compact by condition 4. For any $a>0$, the open radius-$a/3$ balls
centered at points of that compact closure have a finite subcover.
Move each of those finitely many centers to a point of $X$ within
distance $a/3$, using density. The resulting radius-$a$ balls centered
in $X$ cover $A$. Thus $A$ is totally bounded in the ambient-center
convention used here.

Finally assume condition 3. We first prove that every bounded subset
$B\subseteq\widehat X$ is totally bounded in the completion. For any
$a>0$, choose for every $y\in B$ a point $q_y\in X$ with
$m(y,q_y)<a/3$. The set $A=\{q_y:y\in B\}$ is bounded, so it has
a finite radius-$a/3$ cover with centers in $X$. The triangle inequality
shows that these centers cover $B$ by radius-$2a/3$ balls. Since $a$
was arbitrary, $B$ is totally bounded.

Let now $B$ be closed and bounded in $\widehat X$. It is complete,
because it is closed in a complete space, and it is totally bounded by
the preceding paragraph. For completeness, we verify that these two
properties make it compact. Suppose an open cover of $B$ has no finite
subcover. Starting from $C_0=B$, cover $C_{n-1}$ by finitely many
closed balls of radius $2^{-n}$. At least one intersection of such a
ball with $C_{n-1}$ still has no finite subcover from the given open
cover; otherwise the finitely many finite subcovers would cover
$C_{n-1}$. Choose that intersection as $C_n$.

The sets $C_n$ are nonempty, closed in $B$, nested, and have diameter
at most $2^{1-n}$. Choose $x_n\in C_n$. The diameter bound and nesting
make $(x_n)$ Cauchy. Completeness gives a limit $x\in B$, and closedness
gives $x\in C_n$ for every $n$. Some member of the original open cover
contains a relative neighborhood of $x$ in $B$. For sufficiently large
$n$, the diameter bound forces all of $C_n$ into that neighborhood,
contradicting the choice that $C_n$ has no finite subcover. Every open
cover therefore has a finite subcover, so $B$ is compact. This proves
properness of $\widehat X$ and finishes the equivalence.
\end{proof}

\subsection{Examples and interpretation}

\begin{remark}[Example: the rational line]
\label{mt:rationals}
On $X=\mathbb Q$ with ordinary distance, every bounded subset is totally
bounded: subdivide a bounded containing interval into finitely many
intervals of any prescribed small length and choose rational centers.
Thus every countable class of unbounded rational targets satisfies
$\mathsf{NU}_m$. The completion is $\mathbb R$, which is proper.
The original rational space is not proper. For instance,
$[0,2]\cap\mathbb Q$ is closed and bounded in $\mathbb Q$, but a sequence
of rationals converging in $\mathbb R$ to $\sqrt2$ has no convergent
subsequence in $\mathbb Q$, and so this set is not compact.
This distinguishes proper completion from properness of $X$ itself.
\end{remark}

\begin{remark}[Example: a bounded obstruction in an unbounded space]
\label{mt:badspace}
Let $X=\N\times\N$ and
\[
 m\bigl((i,j),(i',j')\bigr)=|i-i'|+\mathbf1_{j\ne j'}.
\]
This is an unbounded metric. The row $A=\{0\}\times\N$ is bounded,
but its distinct points are at distance one. It has no finite
radius-$1/3$ cover. Let
$T_0=\{(i,0):i\ge1\}$ and $T_1=\{(i,1):i\ge1\}$.
The targets $A\cup T_0$ and $A\cup T_1$ are unbounded and have
intersection $A$. They fail nonuniform generation at input radius
$1/3$ and output radius $1$: arbitrarily many distinct row points cross
any proposed input threshold, while one observed row point puts all of
their common intersection within the forbidden output neighborhood.
\end{remark}

\begin{table}[tbp]
\centering
\small
\begin{tabular}{@{}p{.24\linewidth}p{.28\linewidth}p{.39\linewidth}@{}}
\toprule
Scope & Generation property & Exact geometric condition\\
\midrule
One finite target class & Uniform, equivalently nonuniform, at all scales &
Every nonempty finite target intersection is unbounded or totally bounded.\\[5pt]
One countable target class & Nonuniform at all scales &
The same finite-intersection condition.\\[5pt]
Every countable class of unbounded targets & Nonuniform at all scales &
Every bounded subset of the universe is totally bounded; equivalently,
its completion is proper.\\
\bottomrule
\end{tabular}
\caption{Geometric criteria for the semantic generation model of
Definition~\ref{mt:model}. All targets are unbounded.}
\label{mt:table}
\end{table}

The two-target obstruction explains why this is a geometric classification.
Once two targets share a bounded region containing arbitrarily much
small-scale evidence, large input thresholds cannot force a geometrically
novel common output. Conversely, when every bounded finite intersection
is totally bounded, sufficiently dispersed evidence forces an unbounded
common core among the relevant finite prefix of targets. That core always
has room for a novel output.


\section{Contrastive ambiguity and exact list learning}
\label{ca:chapter}
\subsection{Motivation and source papers}
A contrastive observation says that exactly one of two points belongs to the target, but hides which one. How much uncertainty can survive an entire sequence of such observations? If a short list of hypotheses is enough, does that mean the class splits into equally few ordinary identification problems?

These questions combine the list-learning objective of \citet{P14} with the contrastive observation model of \citet{P28}. The latter supplies the common-crossing geometry; the former motivates the comparison between list size and decomposition into identifiable subclasses. The findings below realize arbitrary finite-face ambiguity patterns on countable vertex sets, compute two exact list widths, and identify subclass decomposition with hypergraph coloring. These formulas concern the constructed families, rather than all contrastive classes.

The central example is particularly simple: a class can require only two eventual guesses, yet need arbitrarily many subclasses if each subclass must be identifiable with one guess. Its difficulty is entirely due to contrastive identification: ordinary full-positive-text identification and fresh generation remain easy.

\subsection{Observation and success conventions}
The distinction between the input promises matters throughout this section. A \emph{full positive text} for a support $L$ has range exactly $L$. A \emph{partial positive text} has infinite range contained in $L$, and need not cover $L$. A \emph{contrastive presentation} consists of unordered pairs of distinct points, each pair having exactly one endpoint in $L$, and its endpoints must cover every point of $L$. A deterministic $k$-list learner returns a finite set of at most $k$ target names on every finite history. The empty list is permitted. Lists name the original supports; on histories outside the stated promise a learner may return the empty list. Success means that the true support belongs to every sufficiently late output list; the whole list need not stabilize. All classes are extensional, so different names below denote different supports.

\subsection{Every countable simplicial complex is realizable}
Let $V$ be a nonempty countable vertex set and let $K$ be an abstract simplicial complex of finite subsets of $V$: it contains $\varnothing$ and every singleton, and is closed under taking subsets. Its faces form a countable set. Take disjoint infinite blocks
\[
X=\bigsqcup_{\sigma\in K}B_\sigma,
\qquad B_\sigma=\{(\sigma,n):n\in\N\},
\qquad C=B_\varnothing.
\]
For each $v\in V$ define
\[
N_v=\bigcup_{\substack{\sigma\in K\\v\in\sigma}}B_\sigma,
\qquad h_v=X\setminus N_v,
\qquad \mathcal H_K=\{h_v:v\in V\}.
\]
Every $h_v$ contains the infinite common positive block $C$ and excludes $B_{\{v\}}$. The supports are distinct and incomparable: if $u\ne v$, then $B_{\{u\}}\subseteq h_v\setminus h_u$, and conversely.

For a support $h$ write
\[
\Delta(h)=\{\{x,y\}:x\ne y,\ \text{exactly one of }x,y\text{ lies in }h\},
\qquad
\Gamma(F)=\bigcap_{h\in F}\Delta(h).
\]
A finite nonempty family $F$ is \emph{jointly presentable} if there is one contrastive presentation valid for every member. These are the nonempty faces of the confusability complex of \citet{P28}.

\begin{theorem}[Exact complex realization]\label{thm:ca:realization}
For every finite nonempty $F\subseteq V$, the family $\{h_v:v\in F\}$ is jointly presentable if and only if $F\in K$. Thus every countable abstract simplicial complex is exactly the finite-face contrastive confusability complex of a countable antichain of infinite proper supports with a common infinite positive core.
\end{theorem}
\begin{proof}
Suppose first that a shared presentation exists. Fix $c\in C$. Coverage forces $c$ to appear in an edge $\{c,y\}$. Since $c$ is positive for every target, $y$ must be negative for every $h_v$ with $v\in F$. If $y\in B_\sigma$, this means $F\subseteq\sigma$. Downward closure gives $F\in K$.

Conversely, let $F\in K$ be nonempty. We will cover $U=\bigcup_{v\in F}h_v$ by common crossing edges. Given $x\in U$, let $x\in B_\sigma$ and put $\tau=F\setminus\sigma$. The fact that $x$ is positive for at least one target says $\tau\ne\varnothing$. Since $\tau\subseteq F$, it is a face. Choose $y\in B_\tau$. For every $v\in F$, the point $x$ is negative exactly when $v\in\sigma$, and $y$ is negative exactly when $v\in F\setminus\sigma$. These conditions are complementary on $F$. The blocks are distinct: $\tau$ is nonempty and disjoint from $\sigma$. Therefore $\{x,y\}$ is an edge crossing every target in the family. Enumerate the countable set $U$ and emit one such edge for each point. Every positive point of every target is covered, as required.
\end{proof}

\begin{remark}
Theorem~\ref{thm:ca:realization} controls higher-order faces, not just pairs. A complex may contain every pair of a vertex set while omitting any chosen higher-dimensional faces, subject only to downward closure. Such nonflag behavior is relevant because shared contrastive presentations have genuinely higher-order obstructions.
\end{remark}

\subsection{Three distinct complexity measures}
Define the rank and two-face union bound by
\[
r(K)=\sup_{\sigma\in K}|\sigma|,
\qquad
b(K)=\sup_{\sigma,\tau\in K}|\sigma\cup\tau|.
\]
Both take values in $\N\cup\{\infty\}$. Downward closure implies
\begin{equation}\label{eq:ca:symdiff}
b(K)=\sup_{\sigma,\tau\in K}|\sigma\mathbin{\triangle}\tau|.
\end{equation}
Indeed, a symmetric difference is contained in the union. In the other direction replace $\tau$ by the face $\tau\setminus\sigma$; its symmetric difference with $\sigma$ is $\sigma\cup\tau$. Consequently $r(K)\le b(K)\le2r(K)$ whenever $r(K)$ is finite.

A \emph{nonuniform distinct-edge $k$-list guarantee} uses one learner and a finite threshold $d_v$ for each target: every finite crossing history for $h_v$ with at least $d_v$ distinct unordered edges must return a list containing $h_v$. A uniform guarantee uses one threshold for the entire class. These conditions quantify over finite crossing histories, as do the distinct-edge generation predicates of \citet{P28}. They are stronger than eventual correctness on complete presentations.

\begin{theorem}[Exact eventual and threshold widths]\label{thm:ca:widths}
For every positive integer $k$:
\begin{enumerate}
\item $\mathcal H_K$ is eventually contrastive $k$-list identifiable if and only if $r(K)\le k$. When this holds, the whole list can stabilize.
\item $\mathcal H_K$ has a nonuniform distinct-edge $k$-list guarantee if and only if it has a uniform one, and these hold if and only if $b(K)\le k$. In the positive case, one edge suffices uniformly.
\end{enumerate}
Moreover $\mathcal H_K$ is identifiable from full positive texts and is uniformly contrastively generatable from time zero, with optional avoidance of its own earlier outputs.
\end{theorem}
\begin{proof}
For the first lower bound, any face with $k+1$ vertices supplies a presentation shared by $k+1$ distinct targets, by Theorem~\ref{thm:ca:realization}. Taking the maximum of their success times forces all names into one eventual list. If $r(K)>k$, downward closure supplies such a face.

For the first upper bound, wait until an observed edge involves $C$. Full positive-side coverage guarantees that this happens on every valid presentation. The other endpoint belongs to some $B_\sigma$ with $\sigma\ne\varnothing$. This edge crosses exactly the targets named by $\sigma$, since the endpoint in $C$ is positive for all targets. Store and output that fixed set of names forever. It contains the true target and has size at most $r(K)$. Before that event return a fixed singleton. This also proves whole-list stabilization.

For the threshold upper bound, an edge between $B_\sigma$ and $B_\tau$ crosses $h_v$ exactly when $v\in\sigma\mathbin{\triangle}\tau$. Thus the first edge reveals a fixed candidate set of size at most $b(K)$. Output and retain that set; on a valid history it is nonempty and contains the target.

For necessity, suppose $D=\sigma\mathbin{\triangle}\tau$ has more than $k$ vertices. The distinct infinite blocks $B_\sigma$ and $B_\tau$ contain arbitrarily many distinct unordered edges between them, all of which cross every target named by $D$. If target-dependent thresholds existed, take at least $\max_{v\in D}d_v$ distinct such edges. The resulting one history requires every name in $D$ simultaneously, a contradiction. Each finite history used here extends separately to a full presentation of every target in $D$: append an enumeration of that target paired with a fixed negative point. A shared complete presentation for $D$ is not needed for this finite-prefix argument.

For full-positive-text identification enumerate $V$ and select the least target consistent with the observed sample. On a full text for $h_v$, every earlier wrong $h_u$ is eventually eliminated by a point of $B_{\{u\}}\subseteq h_v\setminus h_u$. There are only finitely many earlier names. Finally, a generator may always choose an element of the infinite common block $C$ outside the finitely many observed endpoints and, if desired, outside all its previous outputs. This succeeds on every history.
\end{proof}

\begin{corollary}[Sharp factor two]\label{cor:ca:two}
For every integer $r\ge1$ and every integer $b$ with $r\le b\le2r$, there is a finite class with exact eventual contrastive list width $r$ and exact nonuniform and uniform distinct-edge list width $b$. It is full-positive-text identifiable and uniformly contrastively generatable from time zero. All supports can be finite unions of residue classes of $\N$, hence unary regular languages.
\end{corollary}
\begin{proof}
Take $r$-element sets $A$ and $B$ whose union has size $b$, allowing $A=B$ when $b=r$. Set $V=A\cup B$, and let $K$ consist of all subsets of either set. Its rank is $r$. Every union of two faces lies in the union of the facets, and the two facets attain size $b$, so its two-face union bound is $b$. Apply Theorem~\ref{thm:ca:widths}. Since $K$ is finite, replace its face blocks by distinct residue classes modulo $|K|$. The same membership and crossing calculations hold.
\end{proof}

\subsection{Hypergraph decomposition is a third invariant}
For $j\ge1$, let $\mathcal E_{j+1}(K)$ be the $(j+1)$-uniform hypergraph on $V$ whose edges are the $(j+1)$-element faces of $K$. A weak coloring has no monochromatic hyperedge.

\begin{theorem}[Exact subclass decomposition]\label{thm:ca:coloring}
The minimum number of eventually contrastive $j$-list-identifiable subclasses covering $\mathcal H_K$ equals the weak chromatic number of $\mathcal E_{j+1}(K)$. In particular, for $j=1$ any countable graph can be realized as the obstruction to decomposing the class into contrastively identifiable subclasses.
\end{theorem}
\begin{proof}
For $W\subseteq V$, the subclass $\{h_v:v\in W\}$ is eventually contrastive $j$-list identifiable exactly when every face of $K$ contained in $W$ has size at most $j$. Necessity follows from the shared-presentation lower bound. For sufficiency use the first-$C$ learner from Theorem~\ref{thm:ca:widths}, but return $\sigma\cap W$ after an edge between $C$ and $B_\sigma$. This is a face contained in $W$, so it has size at most $j$, and it contains the target. Thus every color class gives an admissible subclass in a cover. Conversely, from any finite or countable cover assign to each vertex the least covering index that contains it. A monochromatic $(j+1)$-face would lie in one admissible subclass, contradicting the characterization. Countably many singleton subclasses always provide a cover, so these possibilities exhaust the minimum.
\end{proof}

For example, for $m\ge2$, take all singletons and all pairs of an $m$-vertex complete graph, with no larger faces. The eventual list width is two, but covering by single-identifiable subclasses requires $m$ pieces. The countably infinite complete graph makes the required number countably infinite. This disproves the direct contrastive analog of the full-positive-text $k$-way stratification theorem of \citet{P14}, even though the entire constructed class is already identifiable from ordinary full positive texts.


\subsection{What the finding contributes}
The three quantities separate three questions. The rank $r(K)$ measures ambiguity along a complete presentation. The two-face union bound $b(K)$ measures ambiguity after an arbitrarily long finite prefix. The weak chromatic number measures how many simpler problems are needed to cover the class. Their differences are exact and already occur for finite families of unary regular languages.


\section{Sparse holes and the limits of randomization}
\label{sh:chapter}
\subsection{Motivation and source papers}
An error rate that tends to zero can look like successful learning. Generation in the limit asks for more: after some finite time, every realized output must be valid and fresh. This chapter separates those objectives even for targets that contain almost every natural number, and identifies additional information that makes randomization useful.

The main starting point is the randomized diagonalization of \citet{P10}, which supplies a quantile-cutoff idea. The present construction combines that idea with independently chosen residue classes to make the \emph{missing} points sparse: their consecutive gaps tend to infinity. For every prescribed randomized no-feedback generator, it produces a fixed dense target and a fixed increasing presentation on which the generator makes infinitely many errors almost surely.

Nevertheless, a single generator attains vanishing roundwise error probability and almost-sure zero temporal error frequency for every density-one target and each fixed positive input stream. A further positive result uses a known nondecreasing, unbounded, sublinear sparsity envelope to make the error probabilities summable. Under that common envelope, randomization permits almost-sure eventual correctness even though deterministic generation on all full presentations remains impossible.

All three results concern no-feedback generators and realized outputs. The density promise concerns omitted ambient points, while the error frequency concerns rounds. These differ from density objectives for the set of valid strings generated \citep{P07,P23}.

\subsection{Model and target geometry}
Here a randomized no-feedback generator is a family of measurable maps $g(h,\omega)\in\N$, indexed by finite positive histories $h$, on a fixed random-tape probability space $(\Omega,\Pr)$. Arbitrary dependence between rounds through that tape is allowed. An error means failure of target membership or novelty relative to the observed sample. All probability-one eventual statements refer to realized outputs, without the additional group-representation requirements of \citet{P09}.

Call an infinite set $D=\{d_1<d_2<\cdots\}$ \emph{thin} here if $d_{j+1}-d_j\to\infty$. It has upper Banach density zero:
\[
\lim_{n\to\infty}\sup_{m\ge0}\frac{|D\cap[m,m+n)|}{n}=0.
\]
Indeed, beyond finitely many exceptional points each consecutive gap is at least any prescribed $q$, so every interval contains at most a fixed exceptional count plus $n/q+1$ points. Its complement consequently has one-sided lower Banach density one.

\subsection{A fixed thin-complement target forces almost-sure infinitely many errors}

\begin{theorem}\label{sh:element-diagonal}
For every measurable randomized no-feedback element generator on $\N$, there exist an infinite thin set $D$ and the increasing exact presentation of $L=\N\setminus D$ such that the generator makes infinitely many errors with probability one. Both the target and its presentation are fixed before the random tape is realized.
\end{theorem}
\begin{proof}
Introduce independent construction randomness $R_k$, uniform on $\{0,\ldots,k-1\}$ for $k\ge2$, independent also of the algorithm tape $\omega$. Initialize $B_1=0$ with no holes. Given the previous construction residues, put $M_k=B_{k-1}+k$ and extend the current history increasingly to enumerate every integer at most $M_k$ except holes already selected. Call this history $h_k$.

Conditional on the previous residues, $h_k$ is a fixed finite history independent of $\omega$. Since $g(h_k,\omega)$ is natural-valued, choose the least integer $B_k\ge M_k+k$ satisfying
\[
\Pr_\omega\{g(h_k,\omega)>B_k\}\le2^{-k}.
\]
This cutoff depends on the generator's law and previous construction residues, but not on its realized tape. Reveal the independent residue $R_k$ and declare
\[
D_k=\{n:M_k<n\le B_k,\ n\equiv R_k\pmod k\}
\]
to be additional holes. Let $D=\bigcup_{k\ge2}D_k$.

Every new hole is larger than every currently observed point. The histories are nested increasing lists and their union enumerates exactly $\N\setminus D$. Their lengths tend to infinity because $M_k>B_{k-1}$. Each $D_k$ is nonempty, since its defining interval has length at least $k$. Within $D_k$ the gaps equal $k$, and across successive hole blocks the gap is at least $k+1$. Since each block is finite, $D$ is infinite and thin for every realization of the construction randomness.

Write $Y_k=g(h_k,\omega)$ and $E_k=\{Y_k\le B_k\}$. Independence of the past construction residues from $\omega$ gives
\[
\Pr_{\omega,R}(E_k^c)\le2^{-k}.
\]
Thus only finitely many $E_k$ fail jointly almost surely. Define an auxiliary event by requesting exactly one residue:
\[
Z_k=
\begin{cases}
\{R_k=Y_k\bmod k\},& M_k<Y_k\le B_k,\\
\{R_k=0\},&\text{otherwise}.
\end{cases}
\]
Conditional on the entire algorithm tape and the previous construction residues, the requested residue is fixed while $R_k$ remains uniform. Consequently
\[
\Pr(Z_k\mid\omega,R_2,\ldots,R_{k-1})=1/k.
\]
By iterated conditioning, for $2\le K\le N$,
\[
\Pr\left(\bigcap_{k=K}^N Z_k^c\right)
=\prod_{k=K}^N(1-1/k)=\frac{K-1}{N}.
\]
Letting $N\to\infty$ and taking a countable union over $K$ proves that $Z_k$ occurs infinitely often almost surely. This does not assume that actual generator errors are independent.

On $Z_k\cap E_k$, an actual error occurs at $h_k$. If $Y_k\le M_k$, it is either already observed or an old hole. If $M_k<Y_k\le B_k$, the matching residue makes it a new hole in $D_k$. Since $E_k$ eventually always holds, there are infinitely many actual errors jointly almost surely.

Fubini's theorem now fixes one realization of all construction residues for which the error event has probability one over $\omega$. It gives a deterministic $D$ and a deterministic increasing exact target presentation, proving the assertion. All constructions are measurable: histories and least cutoffs are functions on the countable tree of finite residue prefixes, and membership of each integer is decided after finitely many increasing cutoffs.
\end{proof}

A bound on the generator's output tail, followed by a fresh independent residue, supplies the required probabilistic diagonalization. It is specifically a \emph{no-feedback} proof: with target-dependent external replies, the observed transcript could depend on the random tape and on later-selected holes, requiring a separate independence argument.

\subsection{Vanishing error probability and frequency remain possible}

\begin{theorem}\label{sh:vanishing-error}
There is a randomized no-feedback generator that, for every density-one target and every fixed positive input stream, is always novel, has error probability tending to zero, and has almost-sure temporal error frequency zero.
\end{theorem}
\begin{proof}
At positive history length $t$, independently choose a uniform point from
\[
A_t=\{0,\ldots,t^2\}\setminus S_t.
\]
This set is nonempty because $|S_t|\le t<t^2+1$, and novelty holds identically. At the empty history choose zero, also unseen. For $D=\N\setminus L$,
\[
p_t=\Pr\{g(h_t)\notin L\}
\le\frac{|D\cap[0,t^2]|}{t^2+1-t}\longrightarrow0,
\]
since $D$ has density zero. No full-presentation promise is needed.

For a fixed stream, the error indicators $X_t$ are independent. Put $Z_T=\sum_{t\le T}X_t$. We have $\operatorname{Var}(Z_T)\le T$, while $\mathbb E Z_T/T\to0$ by averaging the convergent sequence $(p_t)$. Chebyshev's inequality at $T=m^2$ gives
\[
 \Pr\{|Z_{m^2}-\mathbb E Z_{m^2}|>\delta m^2\}
 \le\frac1{\delta^2m^2}.
\]
For each rational $\delta>0$ this bound is summable. Borel--Cantelli, followed by a countable intersection over these tolerances, gives $Z_{m^2}/m^2\to0$ almost surely. If $m^2\le T<(m+1)^2$, monotonicity gives
\[
 0\le\frac{Z_T}{T}\le\frac{Z_{(m+1)^2}}{m^2}
 =\frac{(m+1)^2}{m^2}\frac{Z_{(m+1)^2}}{(m+1)^2}.
\]
The right-hand side tends to zero, proving
\[
\frac1T\sum_{t\le T}X_t\longrightarrow0\quad\text{almost surely}.
\]
\end{proof}

Apply Theorem~\ref{sh:element-diagonal} to this generator. On some fixed thin-complement target and its increasing presentation, it makes infinitely many errors almost surely, while their temporal frequency tends to zero almost surely and each round's error probability tends to zero. Thus even this strong geometric target promise does not turn vanishing error into almost-sure eventual correctness.

\subsection{A known sparsity envelope separates randomized and deterministic generation}

Fix a nondecreasing unbounded function $a:\N\to\N$ with $a(n)/n\to0$, and define
\[
\mathcal H_a=\{\N\setminus D:\forall n\ge1,\ |D\cap[0,n)|\le a(n)\}.
\]
For example, take $a(n)=\lceil\log_2(n+1)\rceil$. The envelope is known and common to the class; every target has density one.


\begin{lemma}[Deterministic generation requires a countable inner cover]\label{sh:inner}
Let $\mathcal H$ be any class of infinite subsets of $\N$. If one deterministic no-feedback generator eventually produces valid, previously unobserved points on every full presentation of every target, then there is a countable family $(A_i)$ of infinite subsets of $\N$ such that every $L\in\mathcal H$ contains at least one $A_i$.
\end{lemma}
\begin{proof}
Fix a target $L$. There is a finite positive history $h$ such that the generator succeeds after every finite positive extension of $h$ inside $L$, including the empty extension. Otherwise, start from any finite history over $L$. At each stage, first append the next point of a fixed enumeration of $L$, and then append a finite positive extension at whose end the output is invalid or already observed. The bad extension may be empty, but the enumeration step ensures that the history length increases. The resulting full presentation has infinitely many errors, a contradiction.

For every finite history $h$ over $\N$, run the following deterministic self-feeding process. Query the generator on $h$, append its output, query again, and continue appending the output as the next observation. This process and its range depend only on $h$ and the generator. Retain its range $A_h$ whenever that range is infinite. There are only countably many finite histories, so the retained family is countable. If $h$ is the locking history just constructed for $L$, induction shows that every appended output is in $L$ and outside all preceding observations. The process then has infinite range contained in $L$. Thus every target contains a retained $A_h$.
\end{proof}

\begin{theorem}[Strict deterministic/randomized separation]\label{sh:envelope-separation}
The class $\mathcal H_a$ has no deterministic no-feedback limit element generator under the usual requirement of success on every full target presentation. Nevertheless it has a randomized no-feedback generator of globally infinite sets that, for each target, is almost surely eventually target-valid and fresh, simultaneously on every presentation of that target. For $T\ge1$, the probability of validity and freshness at all times $t\ge T$ is at least $1-2^{1-T}$, uniformly over targets. In particular, randomized eventual element generation succeeds on the class.
\end{theorem}
\begin{proof}
For deterministic impossibility, let $A_i$ be any proposed sequence of infinite inner-cover sets. Revisit each index infinitely often. At stage $m\ge1$, choose a new increasing $d_m\in A_i$ so large that $a(d_m+1)\ge m$, using monotonicity and unboundedness of $a$. For $D=\{d_m:m\ge1\}$, if $[0,n)$ contains $m\ge1$ selected points then $n\ge d_m+1$, so $a(n)\ge m$; the zero-count case is automatic. Thus $D$ satisfies the envelope and intersects every $A_i$ infinitely. Its complement contains none. There is no countable inner cover, so Lemma~\ref{sh:inner} excludes deterministic generation. This necessity concerns all full presentations; the inner-cover argument alone does not establish an increasing-presentation-only lower bound.

For the randomized positive result, choose disjoint intervals $I_j=[N_j,2N_j)$ with $N_j\ge1$ and
\[
\frac{a(2N_j)}{N_j}\le2^{-j},\qquad j\ge1.
\]
Such intervals exist because the ratio tends to zero; choose successive $N_j$ greater than twice their predecessors. Independently select one uniform point $Z_j\in I_j$. At positive history length $t$, emit
\[
Q_t=\{Z_j:j\ge t\}\setminus S_t.
\]
At the empty history output $\{Z_j:j\ge1\}$. The selected points are distinct, so $Q_t$ is infinite for every random tape and every finite history. Removing the sample makes it fresh, including at unrealizable histories.

For each fixed target complement $D$ obeying the envelope,
\[
\Pr(Z_j\in D)\le\frac{a(2N_j)}{N_j}\le2^{-j}.
\]
With probability at least $1-\sum_{j\ge T}2^{-j}=1-2^{1-T}$, no selected point of index at least $T$ is a hole. On that event every $Q_t$ for $t\ge T$ lies in the target and is fresh, for every presentation simultaneously. Borel--Cantelli also gives eventual safety almost surely. Selecting the least element of each infinite $Q_t$ yields the element-valued conclusion.
\end{proof}

The probability-one set of successful tapes may depend on the target. A target chosen after observing the realized tape is a different adversarial model. The infinite output is a semantic random set, not a claim of a finite explicit list encoding. There is no contradiction with Theorem~\ref{sh:element-diagonal}: its diagonal holes need not satisfy any prescribed envelope. A common known density rate permits a summable error budget, whereas the class allowing all unknown rates does not.

\subsection{Comparison of the guarantees}

Appendix Proposition A.2 of \citet{P10} already proves randomized impossibility for the class of all infinite targets, with only finitely many correct outputs almost surely. Its quantile construction makes the target sparse. The new diagonal instead forces infinitely many errors on the complement of a thin set, using quantile bounds together with independent residues. This restricts the target geometry substantially while weakening the error conclusion appropriately; Theorem~\ref{sh:vanishing-error} shows why eventually-all errors cannot be forced for every generator within the density-one class. The refinement concerns the geometry of the admissible targets.

The density objectives in \citet{P07,P23} concern the breadth of generated valid sets. Here, upper Banach density zero is a promise on the omitted ambient points, and the positive theorem controls temporal error frequency. The classes used in the impossibility and separation theorems are uncountable. Thus the conclusions do not contradict positive generation results for arbitrary countable language classes.

\clearpage
\section*{Comparison tables}
\addcontentsline{toc}{section}{Comparison tables}
Tables~\ref{tab:contributions}--\ref{tab:evidence} summarize the mathematical results, their assumptions, and the scope of available verification. The list below also includes the two tables appearing in the main text.
\begingroup
\renewcommand{\listtablename}{List of tables}
\listoftables
\endgroup
\vspace{5pt}
\begin{table}[H]
\centering\small
\caption{The five findings and their source relationships.}\label{tab:contributions}
\begin{tabularx}{\textwidth}{@{}>{\raggedright\arraybackslash}p{.17\textwidth}>{\raggedright\arraybackslash}p{.20\textwidth}Y@{}}
\toprule
\textbf{Finding} & \textbf{Starting point} & \textbf{Additional result and category}\\
\midrule
Overlapping groups & \citet{P09} & Completed-cell characterization, exact two-block discrepancy formula, and sharp square-root rate. Structural refinement and sharp quantitative theorem.\\
\addlinespace
Exact mistake coefficient & \citet{P09,P29} & An explicit coefficient for the minimax logarithmic law, with exact examples and evaluation hardness. Two-paper synthesis, minimax characterization, and complexity.\\
\addlinespace
Metric generation & \citet{P21} & Finite-intersection classification and equivalence to properness of the completion. Geometric characterization informed by a correction to covering conventions.\\
\addlinespace
Contrastive ambiguity & \citet{P14,P28} & Realization of arbitrary ambiguity complexes, exact list widths, and hypergraph-coloring decomposition. Two-paper synthesis and combinatorial separation.\\
\addlinespace
Sparse holes & \citet{P10}; see also \citet{P07,P23} & Dense-target impossibility, vanishing-error guarantees, and strict randomization advantage under a known rate. Probabilistic refinement and separation.\\
\bottomrule
\end{tabularx}
\end{table}
\clearpage
\begin{table}[H]
\centering\small
\caption{Model assumptions that determine the meaning of each theorem.}\label{tab:contracts}
\begin{tabularx}{\textwidth}{@{}p{.19\textwidth}Y Y@{}}
\toprule
\textbf{Finding} & \textbf{Input and structure} & \textbf{Output and success}\\
\midrule
Overlapping groups & One countably infinite target; finite distinct samples; finite dual VC dimension. Sharp rate uses two dual-VC-one blocks and linear completed-cell growth. & A probability law on unobserved target points; discrepancy is worst error over group masses. General infima need not be attained.\\
\addlinespace
Exact mistake coefficient & Fixed finite target class and finite groups; distinct positive observations; target fixed before private randomness; observation order may adapt. & Representation at every round. A draw outside the target or inside the observed sample costs one. Constants in the asymptotic law depend on the fixed instance.\\
\addlinespace
Metric generation & Countable metric universe and countable classes of unbounded targets; positive input and output scales; thresholds measured by ambient covering numbers. & A target point farther than the output scale from the observed sample. Threshold may depend on target and scales, but not on the particular stream.\\
\addlinespace
Contrastive ambiguity & Unordered pairs with exactly one positive endpoint; complete presentations cover the positive side. Threshold results count distinct unordered edges. & A list containing the true original support eventually, or after the specified threshold. Exact formulas apply to the constructed classes.\\
\addlinespace
Sparse holes & Natural-number universe; no feedback; arbitrary measurable private tape. Bad target and increasing full text are fixed before the tape is realized. & Realized output is valid and unobserved. The known-envelope result permits target-dependent probability-one sets of successful tapes.\\
\bottomrule
\end{tabularx}
\end{table}
\vspace{6pt}
\begin{table}[H]
\centering\small
\caption{Which conclusions are different, even when they sound similar.}\label{tab:distinctions}
\begin{tabularx}{\textwidth}{@{}p{.27\textwidth}Y@{}}
\toprule
\textbf{Distinction} & \textbf{Why it matters}\\
\midrule
Small discrepancy vs. an optimizer & Approaching an infimum does not imply some distribution attains it. The proofs retain slack where needed.\\
\addlinespace
Freshness as constraint vs. mistake & Section~\ref{gm:chapter} charges outputs already present in the input sample; this makes immediate representation feasible in depleted groups.\\
\addlinespace
Eventual list width vs. threshold width & Complete contrastive coverage can eventually reveal information that an arbitrarily long finite prefix hides.\\
\addlinespace
Vanishing errors vs. finitely many errors & An error probability and an empirical error fraction can both tend to zero while infinitely many errors occur almost surely.\\
\bottomrule
\end{tabularx}
\end{table}
\clearpage
\begin{table}[H]
\centering\small
\caption{Proof status and provenance of this edition.}\label{tab:evidence}
\begin{tabularx}{\textwidth}{@{}p{.19\textwidth}Y Y@{}}
\toprule
\textbf{Finding} & \textbf{Mathematical proof in these notes} & \textbf{Formal and archival evidence}\\
\midrule
Overlapping groups & Completed-cell iff, exact two-block formula, upper bound, and matching sharpness example. & Later overlap manuscript and audit records. Related finite-profile tools occur in Lean; the complete section is not claimed as formally certified.\\
\addlinespace
Exact mistake coefficient & Matching minimax bounds, explicit examples, and evaluation complexity. & September 18 coefficient manuscript and written audit notes. No Lean certificate for the full logarithmic law.\\
\addlinespace
Metric generation & Finite-intersection iff and equivalence with proper completion, including the negative construction. & September 18 classification manuscript. Covering and fresh-core diagnostics have Lean checks; the full classification is written mathematics.\\
\addlinespace
Contrastive ambiguity & Realization, exact eventual and threshold widths, sharp factor two, and coloring theorem. & F04/F13/F14; \texttt{Round2ContrastiveComplex.lean} has archived successful checks for realization and width equivalences. Coloring is a written consequence.\\
\addlinespace
Sparse holes & Fixed-target diagonalization, vanishing error, inner-cover necessity, and known-envelope separation. & F18/F19 in the September 13 collection, with written AI audit records. Probability layer is not Lean-certified.\\
\bottomrule
\end{tabularx}
\end{table}

\subsection*{Source trail and audit}
This edition consolidates the September 13 \emph{Natural-language proofs}, \emph{Readable finding guide}, and \emph{Ranked discoveries} collection with the September 18 contextual manuscripts \texttt{overlap\_\allowbreak characterization.tex}, \texttt{sharp\_log\_\allowbreak coefficient.tex}, and \texttt{metric\_\allowbreak classification.tex}. The September 19 overlap proof refinements supply the sharp two-block argument. Those files record the wider exploration and additional results not selected here. The section sources retained alongside these notes permit later editing and theorem-level comparison.

A successful archived Lean check certifies the declarations that were checked under their stated definitions. It does not certify neighboring prose, a changed protocol, or the complete collection. The distinctions in Table~\ref{tab:evidence} are therefore part of the notes' mathematical specification.

The five findings were reviewed in a fresh proof-by-proof audit. The review checked the separation and attainment arguments, the minimax dual programs and fixed-target lower bound, the all-scale metric scheduler, the distinction between complete contrastive presentations and finite prefixes, and the measurable fixed-target extraction in the sparse-hole diagonal. It found no material mathematical defect in the stated results. The revision clarifies the positive-supremum attainment claim, integer sample counts, and several presentation and probability conventions.

Supplementary finite checks compared the overlap formula with primal linear programs on 160 instances, checked the minimax constructions through 3,450 linear-program solves, and enumerated all 126 simplicial complexes with one to four vertices. These diagnostics support the review but do not replace the proofs. The accompanying audit records give their precise scope and reproducible code where retained. No new Lean build was performed for this edition; the complete collection is not claimed to be formally certified. The review also does not establish comprehensive literature priority.

\clearpage
\begin{thebibliography}{12}
\addcontentsline{toc}{section}{References}
\raggedright
\bibitem[Charikar et~al.(2025)Charikar, Pabbaraju, and Tewari]{P14} Moses Charikar, Chirag Pabbaraju, and Ambuj Tewari.
\emph{A Characterization of List Language Identification in the Limit}.
arXiv:2511.04103, 2025.
\url{https://arxiv.org/abs/2511.04103}.

\bibitem[Garey et~al.(1976)Garey, Johnson, and Stockmeyer]{GJS1976} Michael R. Garey, David S. Johnson, and Larry Stockmeyer.
\emph{Some simplified NP-complete graph problems}.
Theoretical Computer Science 1(3):237--267, 1976.
\url{https://doi.org/10.1016/0304-3975(76)90059-1}.

\bibitem[Gurvits and Koiran(1997)]{GK} Leonid Gurvits and Pascal Koiran.
\emph{Approximation and learning of convex superpositions}.
Journal of Computer and System Sciences 55(1):161--170, 1997.
\url{https://doi.org/10.1006/jcss.1997.1506}.

\bibitem[Hanneke et~al.(2025)Hanneke, Karbasi, Mehrotra, and Velegkas]{P10} Steve Hanneke, Amin Karbasi, Anay Mehrotra, and Grigoris Velegkas.
\emph{On Union-Closedness of Language Generation}.
arXiv:2506.18642, 2025. See Appendix A.1 for the randomized diagonalization.
\url{https://arxiv.org/abs/2506.18642}.

\bibitem[Kleinberg et~al.(2026)Kleinberg, Peale, and Reingold]{P29} Jon Kleinberg, Charlotte Peale, and Omer Reingold.
\emph{Mistake-Bounded Language Generation}.
arXiv:2605.10809, 2026.
\url{https://arxiv.org/abs/2605.10809}.

\bibitem[Kleinberg and Wei(2025)]{P07} Jon Kleinberg and Fan Wei.
\emph{Density Measures for Language Generation}.
arXiv:2504.14370, 2025.
\url{https://arxiv.org/abs/2504.14370}.

\bibitem[Kleinberg and Wei(2026)]{P23} Jon Kleinberg and Fan Wei.
\emph{Validity, Sparse Holes, and Breadth in Language Generation: Banach Density, Topology, and Geometry}.
arXiv:2604.02385, 2026.
\url{https://arxiv.org/abs/2604.02385}.

\bibitem[Li et~al.(2026a)Li, Raman, and Tewari]{P21} Jiaxun Li, Vinod Raman, and Ambuj Tewari.
\emph{On Generation in Metric Spaces}.
arXiv:2602.07710, 2026a.
\url{https://arxiv.org/abs/2602.07710}.

\bibitem[Li et~al.(2026b)Li, Han, Jiang, and Gao]{P28} Xiaoyu Li, Andi Han, Jiaojiao Jiang, and Junbin Gao.
\emph{Contrastive Identification and Generation in the Limit}.
arXiv:2605.06211, 2026b.
\url{https://arxiv.org/abs/2605.06211}.

\bibitem[Peale et~al.(2025)Peale, Raman, and Reingold]{P09} Charlotte Peale, Vinod Raman, and Omer Reingold.
\emph{Representative Language Generation}.
Proceedings of the 42nd International Conference on Machine Learning, PMLR 267:48518--48541, 2025.
\url{https://proceedings.mlr.press/v267/peale25a.html}.

\bibitem[Vapnik and Chervonenkis(1971)]{VC} V. N. Vapnik and A. Ya. Chervonenkis.
\emph{On the uniform convergence of relative frequencies of events to their probabilities}.
Theory of Probability and Its Applications 16(2):264--280, 1971.
\url{https://doi.org/10.1137/1116025}.
\end{thebibliography}

\end{document}
